% Jacob Neumann

% DOCUMENT CLASS AND PACKAGE USE
    \documentclass[aspectratio=169, handout]{beamer}

    % Establish the colorlambda boolean, to control whether the lambda is solid color (true), or the same as the picture (false)
    \newif\ifcolorlambda
    \colorlambdafalse % DEFAULT: false

    % Use auxcolor for syntax highlighting
    \newif\ifuseaux
    \useauxfalse % DEFAULT: false

    % Color settings
    \useauxtrue

    \newcommand{\auxColor}{2eab9e}     % the color of note boxes and stuff
    \newcommand{\presentColor}{511CE8} % the primary color of the slide borders
    \newcommand{\bgColor}{ebe3ff}      % the color of the background of the slide
    \newcommand{\darkBg}{8b98ad}
    \newcommand{\lambdaColor}{\auxColor}

    \colorlambdatrue

    \usepackage{comment} % comment blocks
    \usepackage{soul} % strikethrough
    \usepackage{listings} % code
    \usepackage{makecell}
    \usepackage{tcolorbox}
    \usepackage{amssymb}% http://ctan.org/pkg/amssymb
    \usepackage{pifont}% http://ctan.org/pkg/pifont
    \usepackage[outline]{contour}
    \usepackage{ stmaryrd }

    \setbeamertemplate{itemize items}[circle]
    % \setbeameroption{show notes on second screen=right}

    \usepackage{lectureSlides}
    %%%%%%%%%%%%%%%%%%%%%%%%%%%%%%%%%%%%%%%%%| <----- Don't make the title any longer than this
    \title{Program Analysis} % TODO
    \subtitle{Writing programs to analyze programs} % TODO
    \date{06 December 2023} % TODO
    \author{Brandon Wu} % TODO


    \def\checkmark{\tikz\fill[green, scale=0.4](0,.35) -- (.25,0) -- (1,.7) -- (.25,.15) -- cycle;}
    \contourlength{.08em}% default is 0.03em
    \newcommand{\cmark}{{\color{green!80!black}\ding{51}}}
    \newcommand{\xmark}{{\color{red}\ding{55}}}


    \newcommand{\dollar}{\mbox{\textdollar}}

    \graphicspath{ {./img/} }


    \colorlet{bgBlue}{hlBlue!130}
    \colorlet{bgOrange}{hlOrange}
    \colorlet{fgBlue}{fgCodeBlue}
    \colorlet{fgGreen}{fgCodeGreen}
    \colorlet{fgRed}{fgCodeRed}
    \colorlet{fgOrange}{fgCodeOrange}

    \colorlet{codeBackground}{background_color}

    \colorlet{hlset}{fgBlue}

    \definecolor{bgPurple}{HTML}{dbc7ff}
    \definecolor{bgRed}{HTML}{fcaea9}
    \definecolor{bgYellow}{HTML}{fffeb5}
    \definecolor{bgGreen}{HTML}{b4ffb3}

    \definecolor{fgOrange}{HTML}{fcae4e}
    \definecolor{fgYellow}{HTML}{ffbe0a}

    \definecolor{unknownPink}{HTML}{e0fff6}

    \usetikzlibrary{positioning}

    \usetikzlibrary{positioning}
    \usetikzlibrary{decorations.pathreplacing,calligraphy}
    \usetikzlibrary{decorations.markings}

    \usetikzlibrary{patterns}
    \usetikzlibrary {calc}
    \tikzset{
      every path/.style={line width=0.25mm},
      tok/.style={
        font=\large,
        align=center,
        minimum height=3cm,
      },
      belowtok/.style={
        font=\large,
        align=center,
        minimum height=3cm,
      },
      astNode/.style={
        draw,
        inner sep=2pt,
        line width=0.4mm,
        minimum size=0.6cm,
        font=\large,
        align=center,
      },
      hex/.style={
        astNode,
        signal,
        signal to=east and west,
        signal pointer angle=130,
        fill=bgPurple,
      },
      box/.style={
        astNode,
        rectangle,
        fill=green!20!white,
      },
      bbox/.style={
        astNode,
        rectangle,
        fill=bgBlue,
      },
      circ/.style={
        astNode,
        circle,
        minimum size=0.7cm,
        fill=bgYellow
      },
      bhex/.style={
        astNode,
        signal,
        inner sep=0.5pt,
        signal to=east and west,
        signal pointer angle=130,
        fill=bgBlue,
        minimum width=0.8cm,
      },
      unknown/.style={hex, draw=black, preaction={fill, unknownPink}, pattern=north east lines, pattern color=gray},
      unknown2/.style={unknown, preaction={fill, bgYellow}},
      highlight/.style={draw=fgYellow, line width=0.55mm},
      highlight2/.style={draw=magenta, line width=0.55mm},
      highlight3/.style={draw=blue, line width=0.55mm},
      highlight4/.style={draw=orange, line width=0.55mm},
      highlight5/.style={draw=green!70!white, line width=0.55mm},
      hlbg/.style={fill=bgYellow},
      hledge/.style={thick, red},
      between/.style args={#1 and #2}{
        at = ($(#1)!0.5!(#2)$)
      },
    }

    % DONT FORGET TO PUT [fragile] on frames with codeblocks, specs, etc.
        %\begin{frame}[fragile]
        %\begin{codeblock}
        %fun fact 0 = 1
        %  | fact n = n * fact(n-1)
        %\end{codeblock}
        %\end{frame}

    % INCLUDING codefile:
        % 1. In some file under code/NN (where NN is the lecture id num), include:
    %       (* FRAGMENT KK *)
    %           <CONTENT>
    %       (* END KK *)

    %    Remember to not put anything on the same line as the FRAGMENT or END comment, as that won't be included. KK here is some (not-zero-padded) integer. Note that you MUST have fragments 0,1,...,KK-1 defined in this manner in order for fragment KK to be properly extracted.
        %  2. On the slide where you want code fragment K
                % \smlFrag[color]{KK}
        %     where 'color' is some color string (defaults to 'white'. Don't use presentColor.
    %  3. If you want to offset the line numbers (e.g. have them start at line 5 instead of 1), use
                % \smlFragOffset[color]{KK}{5}

\begin{document}

% Make it so ./mkWeb works correctly
\ifweb
    \renewcommand{\pause}{}
\fi

\setbeamertemplate{itemize items}[circle]

% SOLID COLOR TITLE (see SETTINGS.sty)
{
\begin{frame}[plain]
    \colorlambdatrue
    \titlepage
\end{frame}
}

\begin{frame}[fragile]
  \frametitle{Lesson Plan}

  \tableofcontents
\end{frame}

\sectionSlide{1}{The State of Software}

\begin{frame}[fragile]
  \frametitle{Software and the World}

  On August 20th, 2011, Silicon Valley venture capitalist and and
  entrepreneur Marc Andreessen\footnote{Currently a board director for Meta Platforms.} published an essay entitled
  "Software is eating the world".

  \pause
  \vspace{\fill}

  This essay included a lot of business-oriented reasons for why software
  was immensely disrupting each individual economic sector, for reasons of
  ease of use, speed of execution, and reach of influence, among others.

  \pause
  \vspace{\fill}

  Now, more than a decade after this article, it's an incredibly obvious fact that
  software already has eaten the world. You cannot get away from it -- it is
  everywhere, and it is everything.
\end{frame}

\begin{frame}[fragile]
  \frametitle{The State of Software Engineering}

  Part of what makes software engineering a lucrative profession is that there
  is not, and will never be, a shortage of need for software engineers.

  \pause
  \vspace{\fill}

  Regardless of if a company is a recruiting company, a think tank, a massage parlor,
  or a pet food retailer, everyone needs software developers. Every business needs
  a website, every business deals with data, and every business needs a way to
  keep up with every other business, which is doing exactly the same.

  \pause
  \vspace{\fill}

  Unfortunately, not all of them are Jonathan Werrett, and are well-informed on
  the importance of writing safe code.
\end{frame}

\begin{frame}[fragile]
  \frametitle{The State of Software Engineering}

  One theme that crops up in the general practice of software engineering is to
  try to produce as little code as possible, because any human writing any
  amount of code has some probability of producing a bug.

  \pause
  \vspace{\fill}

  The less code we write, the less possibility of writing a bug.

  \pause
  \vspace{\fill}

  So what can we say about the immense volume of code produced by
  the tens of millions of software developers around the world?

  \pause
  \vspace{\fill}

  \textbf{Answer:} It is horribly, immensely buggy, and full of mistakes.
\end{frame}

\begin{frame}[fragile]
  \frametitle{What's in an Error?}

  When you write a mistake in your code, what does it often look like?

  \pause
  \vspace{\fill}

  Maybe you made a typo:

  \begin{center}
    \begin{minipage}[t][1.0in][t]{\textwidth}
      \begin{minipage}{0.2\textwidth}
        \centering
        \vspace{\fill}
        {\huge\xmark}
        \vspace{\fill}
      \end{minipage}
      \vspace{\fill}
      \begin{minipage}{0.75\textwidth}
        {\small
        \begin{codeblock}[rulecolor=\color{red}, framerule=0.3mm, language=python]
          def fact(n):
            if n == 0:
              return 1
            return n * fac(n - 1)
        \end{codeblock}
        }
      \end{minipage}
    \end{minipage}
    \pause
    \begin{minipage}[t][1.0in][t]{\textwidth}
      \begin{minipage}{0.2\textwidth}
        \centering
        \vspace{\fill}
        {\huge\cmark}
        \vspace{\fill}
      \end{minipage}
      \vspace{\fill}
      \begin{minipage}{0.75\textwidth}
        {\small
        \begin{codeblock}[rulecolor=\color{green!80!black}, framerule=0.3mm, language=python]
          def fact(n):
            if n == 0:
              return 1
            return n * `fact`(n - 1)
        \end{codeblock}
        }
      \end{minipage}
    \end{minipage}
  \end{center}
\end{frame}

\begin{frame}[fragile]
  \frametitle{What's in an Error?}

  Or maybe you declared a variable, and then forgot to use it:

  \pause
  \begin{center}
    \begin{minipage}[t][0.9in][t]{\textwidth}
      \begin{minipage}{0.2\textwidth}
        \centering
        \vspace{\fill}
        {\huge\xmark}
        \vspace{\fill}
      \end{minipage}
      \begin{minipage}{0.75\textwidth}
        {\small
          \begin{codeblock}[rulecolor=\color{red}, framerule=0.3mm, language=python]
            fun compoundInterest(n, rate):
              amountToAdd = n * rate
              return n
          \end{codeblock}
          }
        \end{minipage}
    \end{minipage}
    \pause
    \begin{minipage}[t][0.9in][t]{\textwidth}
      \begin{minipage}{0.2\textwidth}
        \centering
        \vspace{\fill}
        {\huge\cmark}
        \vspace{\fill}
      \end{minipage}
      \begin{minipage}{0.75\textwidth}
        {\small
        \begin{codeblock}[rulecolor=\color{green!80!black}, framerule=0.3mm, language=python]
          fun compoundInterest(n, rate):
            amountToAdd = n * rate
            return `n + amountToAdd`
        \end{codeblock}
        }
      \end{minipage}
    \end{minipage}
  \end{center}
\end{frame}

\begin{frame}[fragile]
  \frametitle{What's in an Error?}

  Or maybe you just have a simple type error:

  \pause
  \begin{center}
    \begin{minipage}[t][1.1in][t]{\textwidth}
      \begin{minipage}{0.2\textwidth}
        \centering
        \vspace{\fill}
        {\huge\xmark}
        \vspace{\fill}
      \end{minipage}
      \begin{minipage}{0.75\textwidth}
        {\small
          \begin{codeblock}[rulecolor=\color{red}, framerule=0.3mm, language=python]
            def initialize(data):
              d = {}
              for s in data:
                d[s.strip] = None
              return d
          \end{codeblock}
          }
        \end{minipage}
    \end{minipage}
    \pause
    \begin{minipage}[t][1.1in][t]{\textwidth}
      \begin{minipage}{0.2\textwidth}
        \centering
        \vspace{\fill}
        {\huge\cmark}
        \vspace{\fill}
      \end{minipage}
      \begin{minipage}{0.75\textwidth}
        {\small
        \begin{codeblock}[rulecolor=\color{green!80!black}, framerule=0.3mm, language=python]
            def initialize(data):
              d = {}
              for s in data:
                d[s.strip`()`] = None
              return d
        \end{codeblock}
        }
      \end{minipage}
    \end{minipage}
  \end{center}
\end{frame}

\begin{frame}[fragile]
  \frametitle{A Machine Learning Parable}

  For most of these errors, they can seem innocent enough. Most are able to be
  identified and fixed before code ever sees the light of day.

  \pause
  \vspace{\fill}

  Not all are so lucky, however. Imagine that you are a machine learning engineer.

  \pause
  \vspace{\fill}

  You have spent the past six months working on a state of the art large
  language model\footnote{Henceforth referred to as "NoCapGPT"}, and finally you are ready to put it to the test. You just
  make a few adjustments (mostly adding comments and clarifying names),
  before you run the model and then decide to go on a vacation to France for
  two weeks.

  \pause
  \vspace{\fill}

  When you return from your vacation, you discover that your model failed
  with:
  {\small
  \begin{lstlisting}[style=15150code, numbers=none]
    NameError: name 'modle' is not defined. Did you mean: 'model'?
  \end{lstlisting}
  }

  \pause
  \vspace{\fill}

  This can actually happen.
\end{frame}

\begin{frame}[fragile]
  \frametitle{Explosive Repercussions}

  What's the point? A computer program can be compared to a loaded pistol.

  \pause
  \vspace{\fill}

  To know whether a pistol is defective, we could just shoot it, and see whether
  or not it blows up. This is analogous to detecting an error at \term{runtime},
  when actually running a given program.

  \pause
  \vspace{\fill}

  Or, we could be a little more thorough about it, and inspect the pistol for
  damage, or if someone put .50 AE into a 9mm pistol\footnote{I had to look this
  up. Basically, big ammo goes in small gun.}.

  \pause
  \vspace{\fill}

  Obviously, the solution which ends up with the gun not exploding in our hands is
  the preferable one.
\end{frame}

\begin{frame}[fragile]
  \frametitle{A Better Way}

  This is just another way of saying that we need a better method than just
  running into errors blindly, when running code in real environments. Discovering
  a bug in the Curiosity rover before it takes off is a miracle. Discovering a bug
  in the Curiosity rover \textit{after} it takes off is a failure.

  \pause
  \vspace{\fill}

  This is the basic principle behind \term{shift left}, which means evaluating
  our programs by some criteria as far \textit{before} runtime as possible, or
  farther left, chronologically.

  \pause
  \vspace{\fill}

  To do this, we will need a better way of catching these errors, of analyzing
  whether the programs we write contain code smells, fatal errors, or even
  security vulnerabilities.
\end{frame}



\begin{frame}[fragile]
  \frametitle{The Cost of Mistakes}

  What's the point? \textbf{Programmers make mistakes}.

  \pause
  \vspace{\fill}

  There are many programmers in the world. Programs that make these kinds of
  silly, one-off mistakes happen millions of times in a single day. For every
  mistake which is caught before reaching a sensitive application, there are
  millions more waiting to be uncovered. We need to build tools, compilers, and
  programming languages that can guard against these mistakes because the
  cost of allowing them is too high.

  \pause
  \vspace{\fill}

  Tony Hoare once referred to his invention of the null pointer\footnote{A
  programming idiom which has led to countless mistakes.} as his
  \textit{billion-dollar mistake}. We are talking about fighting a war upon
  which rests not only billions upon billions of dollars across every
  conceivable industry, but upon which rests the security and continued
  operation of our society.
\end{frame}

\begin{frame}[fragile]
  \frametitle{Battling Software}

  How can we prevent these bugs? How do we make sure that,
  for the prodigious, gargantuan, and overflowing deluge of software
  that is pumped out every day, it is as safe and as correct as possible?
  The alternative is a reality that is horrifying to contemplate.
  \footnote{I highly recommend the book \textit{This Is How They Tell Me
  the World Ends: The Cyberweapons Arms Race} by Nicole Perlroth.}

  \pause
  \vspace{\fill}

  Software is eating the world.

  \pause
  \vspace{\fill}

  It's time to bite back.
\end{frame}

\sectionSlide{2}{Program Analysis}

\begin{frame}[fragile]
  \frametitle{What is Program Analysis?}

  \defBox{}{\term{Program analysis} is the art of discovering undesirable
  behaviors in programs, usually by automated, programmatic means.}

  \pause
  \vspace{\fill}

  These undesirable behaviors may include correctness, performance,
  security, and legibility. Ultimately, it encompasses any property
  which is worth testing, of a program.

  \pause
  \vspace{\fill}

  In essence, program analysis entails \textbf{writing programs to
  analyze programs}.

  \pause
  \vspace{\fill}

  \begin{figure}
  \begin{codeblock}[language=python]
    def analyze(program):
      if programIsBad(program):
        dont()
  \end{codeblock}
    \caption{Figure 1: An extremely simplified view of the entire \code{semgrep} repository.}
  \end{figure}
\end{frame}

% \begin{frame}[fragile]
%   \frametitle{Refresher: Interpreted and Compiled Languages}

%   The lifecycle of a program is varied, however. When should this analysis take place?

%   \pause
%   \vspace{\fill}

%   As a refresher, recall the notion of programming languages as \term{interpreted} or \term{compiled}.

%   \pause
%   \vspace{\fill}

%   A program is a text file written in some kind of programming language. Some languages
%   are written in \term{interpreted languages}, meaning that they can be executed (or \textit{run})
%   immediately.

%   \pause
%   \vspace{\fill}

%   Other languages are \term{compiled languages}, meaning that an executable version of
%   the program must first be obtained by running a separate program called the
%   \term{compiler}.
% \end{frame}

% \begin{frame}[fragile]
%   \frametitle{An Americentric Analogy}

%   For an Americentric analogy, you can think of a program in an interpreted language as
%   a recipe written in English, which the average American can follow just by reading.
%   A program in a compiled language is written in French, which, for the average American,
%   must first be translated to English before using.

%   \pause
%   \vspace{\fill}

%   \begin{center}
%     \begin{tikzpicture}
%       \node[hex, minimum size=1cm] (A) {English recipe};
%       \node[hex, right of=A, node distance=2in, minimum size=1cm, fill=red!30!white] (B) {food};

%       \draw[-stealth] (A) -- (B) node[midway, above] {cooking};
%     \end{tikzpicture}

%     \vspace{\fill}

%     \begin{tikzpicture}
%       \node[box, minimum size=1cm, fill=orange!40!white] (A) {French recipe};
%       \node[hex, minimum size=1cm, right of=A, node distance=2.5in] (B) {English recipe};
%       \node[hex, right of=B, node distance=2.5in, minimum size=1cm, fill=red!30!white] (C) {food};

%       \draw[-stealth] (A) -- (B) node[midway, above] {translation};
%       \draw[-stealth] (B) -- (C) node[midway, above] {cooking};
%     \end{tikzpicture}
%   \end{center}
% \end{frame}

% \begin{frame}[fragile]
%   \frametitle{Refresher: Interpreted and Compiled Languages}

%   \begin{center}
%     \begin{tikzpicture}
%       \node[hex, minimum size=1cm] (A) {program};
%       \node[hex, right of=A, node distance=2in, minimum size=1cm, fill=red!30!white] (B) {result};

%       \draw[-stealth] (A) -- (B) node[midway, above] {execution};
%     \end{tikzpicture}

%     \vspace{\fill}

%     \begin{tikzpicture}
%       \node[box, minimum size=1cm, fill=orange!40!white] (A) {program};
%       \node[hex, minimum size=1cm, right of=A, node distance=2.5in] (B) {compiled program};
%       \node[hex, right of=B, node distance=2.5in, minimum size=1cm, fill=red!30!white] (C) {result};

%       \draw[-stealth] (A) -- (B) node[midway, above] {compilation};
%       \draw[-stealth] (B) -- (C) node[midway, above] {execution};
%     \end{tikzpicture}
%   \end{center}
% \end{frame}


\begin{frame}[fragile]
  \frametitle{Flavors of Program Analysis: Dynamic}

  Program analysis generally comes in one of two flavors:

  \pause
  \vspace{\fill}

  \defBox{}{\term{Dynamic program analysis} has to do with figuring
  out program behavior by observing its behavior at run-time.}

  \pause
  \vspace{\fill}

  This can include things like \term{fuzzing}, which involves running
  the program on a wide range of random inputs, \term{profiling}, which
  involves measuring the run-time of a program on some inputs, and even
  the simple act of writing tests.

  \pause
  \vspace{\fill}

  Dynamic program analysis is useful, and goes straight to the source in terms
  of the program's actual behavior, but it is limited in some other ways.
  Notably, if the target program loops forever or takes a really long time, then
  dynamic program analysis will do the same.

  \pause
  \vspace{\fill}

  Another is that this doesn't necessarily stop the code from being written or
  committed\footnote{"uploaded", essentially.} in the first place. We want something
  even farther "left", if possible.

\end{frame}

\begin{frame}[fragile]
  \frametitle{Flavors of Program Analysis: Static}

  \defBox{}{\term{Static program analysis} concerns ascertaining properties
  of programs \textbf{without ever running the program}.}

  \pause
  \vspace{\fill}

  This has the advantage that it can happen before code ever leaves the programmer's
  computer.

  \pause
  \vspace{\fill}

  This will be our focus for today. Specific applications of this analysis include:
  \pause
  \begin{itemize}
    \item \term{static application security testing} (or SAST), which is
    the process of applying static program analysis to code for security purposes \pause
    \item syntax highlighting, which looks at a (possibly incomplete) program
    and tries to color it, as its being written \pause
    \item autoformatting, which looks at a program and tries to make it adhere to
    a certain stylistic convention \pause
    \item type-checking, which looks at a program and ascertains what type its
    constituent parts have (if any)
  \end{itemize}
\end{frame}

\begin{frame}[fragile]
  \frametitle{A Minor Issue}

  This is the mission we have ahead of us. Before we can dive into more
  technical details, however, we have one small issue before us:

  \pause
  \vspace{\fill}

  \textbf{Program analysis is inherently impossible}.

  \pause
  \vspace{\fill}

  {\color{blue}\href{https://en.wikipedia.org/wiki/Rice%27s_theorem}{Rice's
  Theorem}} is a mathematical theorem in computability theory\footnote{The field
  of computability theory is kind of similar to forum posts speculating on what
  would happen in a fight between Batman and Goku. It's just nerds getting
  together and thinking about increasingly more complicated levels of
  impossibility.} which states:
  \begin{quote}
    All non-trivial semantic properties of programs are undecidable.
  \end{quote}

  \pause
  \vspace{\fill}

  In English:
  \begin{quote}
   It is impossible to definitively answer yes or no for any property of a
   program's behavior, in a finite amount of time.
  \end{quote}
\end{frame}

\begin{frame}[fragile]
  \frametitle{The Halting Problem}

  This is a corollary of the
  {\color{blue}\href{https://en.wikipedia.org/wiki/Halting_problem}{Halting
  Problem}}, which states that it is impossible to write a program to tell if a
  program loops forever or not.\footnote{Note that these claims of impossibility
  are \textit{in general}. For instance, I can look at the program \code{def
  loop(): loop()} with my eyes and tell you that it loops forever, but we cannot
  write a program which does that for \textit{every} program.} The reason why
  this is impossible come out of asking a simple question: what should the
  following function return?

  \pause
  \vspace{\fill}

  {\small
  \begin{codeblock}[language=python]
    def foo():
      if halts(foo):
        loop()
      else:
        return
  \end{codeblock}
  }

  \pause
  \vspace{\fill}

  This ends up producing a paradox, in much the same way as the
  {\color{blue}\href{https://en.wikipedia.org/wiki/Liar_paradox}{liar's paradox}}.

  \pause
  \vspace{\fill}

  Because any program can loop forever, this ends up tainting every other
  question that program analysis could answer, meaning that all of them are inherently
  impossible.
\end{frame}

\begin{frame}[fragile]
  \frametitle{Battling Impossibility}

  Well, that sucks. So then, with this knowledge, is this talk over?

  \pause
  \vspace{\fill}

  No, because \textbf{it only sucks if you're a quitter.}

  \pause
  \vspace{\fill}

  Generally, when solving a hard problem, it's usually either a sign that we
  are not working hard enough, or our expectations are too high.
  \pause
  \vspace{\fill}

  Well, in this case, this problem's impossibility is a mathematical truth, so
  it's not a skill issue in terms of our ability to implement. That's not the
  problem here.

  \pause
  \vspace{\fill}

  So let's lower our expectations.

  \pause
  \vspace{\fill}

  \begin{quote}
    It is impossible to definitively answer yes or no for any property of a
    program's behavior, in a finite amount of time.
  \end{quote}

  \pause
  \vspace{\fill}

  The main innovation out of program analysis is -- \textit{it's only impossible if
  you insist on being right all the time}.
\end{frame}

\begin{frame}[fragile]
  \frametitle{Compilers and Program Analysis}

  Generally, we want our tools to be trustworthy. The day that a pacemaker or
  ladder stops working is the day that someone gets hurt. Nobody wants to use
  something which only works \textit{some of the time}.

  \pause
  \vspace{\fill}

  So, for most products, they must work \textit{all the time}. To do otherwise
  would be sure way to head straight for bankruptcy in the free
  marketplace.

  \pause
  \vspace{\fill}

  Program analysis suffers from no such thing. Our goal will be to implement
  analyses which always complete within a finite amount of time, albeit with
  the caveat that sometimes they might be inaccurate or incomplete, or throw their
  hands up and say "I don't know".
\end{frame}

\begin{frame}[fragile]
  \frametitle{Full Employment}

  A funny corollary of this sentiment is something called the
  {\color{blue}\href{https://en.wikipedia.org/wiki/Full-employment_theorem}{full-employment theorem}},
  which essentially states that there will always be jobs in program analysis and
  compiler-writing, i.e. employment is always ensured\footnote{Knock on wood.}.

  \pause
  \vspace{\fill}

  This is because, due to the fact that the task is inherently impossible, it's
  always possible to write a better analysis that covers more cases, or a
  compiler which can produce better binaries. You just need more casework.
\end{frame}

\begin{frame}[fragile]
  \frametitle{Program Analysis via Casework}

  For instance, no one is stopping you from doing this:
  \pause
  \begin{codeblock}[language=python]
    def halts(program: str) -> Optional[bool]:
      if program == "x = 1":
        return true
      elif program == "x = 2":
        return true
      elif program == "def loop:\n loop()":
        return false
  \end{codeblock}

  \pause
  \vspace{\fill}

  Far from the cutting edge, but it works. A team of monkeys at typewriters
  could eventually churn out a more effective \code{halts} function than
  exists anywhere else.

  \pause
  \vspace{\fill}

  You might find this demoralizing. I find it motivating, somehow.
\end{frame}

\begin{frame}[fragile]
  \frametitle{Resolution}

  So, this is the scope of the task in front of us.

  \pause
  \vspace{\fill}

  We have finite resources and finite time to solve a problem which is
  impossible to solve.

  \pause
  \vspace{\fill}

  And still, software is eating the world. The cost of failing is too high.

  \pause
  \vspace{\fill}

  Time to put in some elbow grease and get to work.
\end{frame}

\sectionSlide{3}{Doing the Impossible}

\begin{frame}[fragile]
  \frametitle{Type Checking}

  Consider a more specific example of program analysis, namely that of
  type-checking a program.

  \pause
  \vspace{\fill}

  Consider a function \code{f} which adds two ints together. We might say that
  it has the \term{type} of \code{(int, int) -> int} -- that is, it takes in
  two integers and produces an integer. Automatically figuring this out is the
  problem of type-checking.

  \pause
  \vspace{\fill}

  Let's consider the problem of giving a type to a \code{div}, which
  takes two integers and computes their integer division.
\end{frame}

\begin{frame}[fragile]
  \frametitle{A More Accurate \code{div}}

  We might say that it also has type \code{(int, int) -> int}, because it seems
  to take two \code{int}s and return an integer.

  \pause
  \vspace{\fill}

  There is one caveat, however. If we pass in \code{0} as the second parameter,
  it will crash on us!

  \pause
  \vspace{\fill}

  Well, this is a dynamic error. We want to catch such things before running
  the program. Can we assign \code{div} a type, such that it will prevent us
  from ever running code which passes \code{0} to \code{div}?

  \pause
  \vspace{\fill}

  There are two properties that we would desire of such a type:
  \begin{itemize}
    \item we can ascertain it in a finite amount of time \pause
    \item it says that a usage of \code{div} is ill-typed if and only if
    we pass \code{0}, or a non-\code{int}, to it
  \end{itemize}

  \pause
  \vspace{\fill}

  \textbf{It is impossible to have both of these things at the same time.}
\end{frame}

\begin{frame}[fragile]
  \frametitle{Sweaty Palms}

  \begin{center}
    \includegraphics[scale=0.3]{sweaty}
  \end{center}
\end{frame}

\begin{frame}[fragile]
  \frametitle{Dependent Types and Type A Analysis}

  There is a class of languages which have very sophisticated type systems,
  called \term{dependently typed} languages, where you can actually express
  this.

  \pause
  \vspace{\fill}

  The issue is that type-checking in such languages can take forever,
  meaning we would get the second property, but not the first. We will
  call this a \term{Type A program analysis}.

  \pause
  \vspace{\fill}

  So the other track will be to get the first property, but not the second.
  We will have to accept being wrong sometimes, and either rejecting
  programs which do not divide by zero, or accepting programs which do.
  We will call this \term{Type B program analysis}.
\end{frame}

\begin{frame}[fragile]
  \frametitle{An Analogy for Type B Analysis}

  Consider the following analogy.

  \pause
  \vspace{\fill}

  You are in charge of security at an airport.

  \pause
  \vspace{\fill}

  You are acutely aware of the fact that in the early 2000s, a man tried to set off
  plastic explosives concealed in his shoes, during a transatlantic flight
  from France to Florida.

  \pause
  \vspace{\fill}

  The issue is that you don't have a good way of telling whether an arbitrary
  individual might have explosives in their shoes, or not.

  \pause
  \vspace{\fill}

  So, how do you minimize the chances?

  \pause
  \vspace{\fill}

  This is the story of why everyone needs to take off their shoes in the airport.
\end{frame}

\begin{frame}[fragile]
  \frametitle{A Rough Approximation}

  The parable of this story is that, while it may be difficult or impossible to
  gather an \textit{exact} answer (who has explosives in their shoes), it's easy
  to obtain an \textit{approximative} answer: assume that \textit{everyone} has
  explosives in their shoes.

  \pause
  \vspace{\fill}

  So just scan everybody's shoes. Problem solved.

  \pause
  \vspace{\fill}

  So how can we tell which programs contain an unsafe call to \code{div}? Well,
  if we don't mind being wrong sometimes...

  \pause
  \vspace{\fill}

  \begin{codeblock}[language=python]
    def containsUnsafeDiv(program: str) -> bool:
      return "div" in program
  \end{codeblock}
\end{frame}

\begin{frame}[fragile]
  \frametitle{Approximation in Practice}

  Let's see how it behaves:

  \begin{itemize}
    \item \code{containsUnsafeDiv("x = 2")} returns \code{False}. \cmark \pause
    \item \code{containsUnsafeDiv("div(1, 0)")} returns \code{True}. \cmark \pause
    \item \code{containsUnsafeDiv("div(1, 1)")} returns \code{True}. \xmark \pause
    \item \code{containsUnsafeDiv("// i like division!")} returns \code{True}. \xmark
  \end{itemize}

  \pause
  \vspace{\fill}

  We see that it works... sometimes.

  \pause
  \vspace{\fill}

  But this is precisely what a Type B program analysis purports to do. We accept
  that we might be wrong, sometimes, in exchange for the fact that this test always
  completes in short order.
\end{frame}



\begin{frame}[fragile]
  \frametitle{Tradeoffs}

  Which outcome is more preferred? Well, it turns out the answer is
  "neither of them".

  \pause
  \vspace{\fill}

  We are basically saying that, to be able to reject all programs which might
  divide by zero, we can either accept infinitely looping compile times, or we
  can reject every program which contains the characters \code{div} in its
  source text.

  \pause
  \vspace{\fill}

  Now, with more sophisticated techniques, we can do a little bit better than
  rejecting every program containing \code{div}. But not by much.

  \pause
  \vspace{\fill}

  It turns out, in practice, the right solution will be to simply not
  care so much about the division by zero case. It's not worth the
  trade-offs.
\end{frame}

\begin{frame}[fragile]
  \frametitle{Downscoping and Type C Analysis}

  We said it was impossible to have the virtues of Type A and Type B analysis at the
  same time. That's true, if we fix the problem statement. We might say that
  \term{Type C program analysis} is to both terminate and be correct, but at the cost of
  simplifying the problem we are trying to solve.

  \pause
  \vspace{\fill}

  Type-checking is usually an example of a Type C analysis. So, thus we end up
  with not being able to statically catch division by zero errors. For the
  question of "does this program divide by zero?", we decided the answer is "we
  don't care".

  \pause
  \vspace{\fill}

  What about the question of "does this program divide by a non-integer"?

  \pause
  \vspace{\fill}

  It turns out, this is perfectly solvable in a terminating manner.\footnote{For the
  purposes of this presentation, it's not actually important this problem is solvable.
  Just know there are some program analysis problems which can be solved without
  looping or approximating (but few).} The
  reason why this is OK is that "non-integer" is an approximative query --
  "zero" is specific.
\end{frame}

\begin{frame}[fragile]
  \frametitle{Recap}

  So let's recap for a second:
  \pause
  \begin{itemize}
    \item we would like to answer specific questions about programs, which are
    impossible to do in general. \pause
    \item We need to give up one of guaranteed termination, perfect accuracy,
    or solving that exact problem. Types A, B, and C analysis correspond to
    giving up each of these things, respectively.
  \end{itemize}

  \pause
  \vspace{\fill}

  Specific examples include:
  \begin{itemize}
    \item dependent typechecking is a Type A analysis (can loop forever) \pause
    \item rejecting all programs with \code{div} is a Type B analysis (rejects valid programs) \pause
    \item regular typechecking is a Type C analysis (give up catching divide by zero)
  \end{itemize}

  \pause
  \vspace{\fill}

  For most practical program analysis tools, looping forever is not an option.\footnote{We
  might call this "doubly impossible".} So for our purposes, we are generally
  interested in Type B and Type C analysis.
\end{frame}

\begin{frame}[fragile]
  \frametitle{Flavors of Program Analysis}

  An easy way to remember program analysis is to envision it as three different kinds of people.

  \pause
  \vspace{\fill}

  \begin{minipage}{0.32\textwidth}
    \begin{center}
      \includegraphics[scale=0.2]{oldwoman.jpg}
      Type A analysis is personified by an old person (they might take forever to answer)
    \end{center}
  \end{minipage}
  \pause
  \begin{minipage}{0.32\textwidth}
    \begin{center}
      \includegraphics[scale=0.1]{flatearther.jpg}
  Type B analysis is personified by a flat earther (they sometimes spout nonsense)
    \end{center}
  \end{minipage}
  \pause
  \begin{minipage}{0.32\textwidth}
    \begin{center}
      \includegraphics[scale=0.8]{kid.jpg}
  Type C analysis is personified by a literal child (they can only answer simple questions)
    \end{center}
  \end{minipage}
\end{frame}

 \begin{frame}[fragile]
   \frametitle{On Incorrectness}

   But, we still have significant problems left to answer.

   \pause
   \vspace{\fill}

   Type-checking is only in the Type C category because of decades of work by type theorists
   and language designers, to figure out what buckets of questions can be answered tractably
   by machines, and to what extent.

   \pause
   \vspace{\fill}

   For many other problems, such as code reachability, vulnerability to outside
   attackers, and unsafe behavior, we have no such guarantees. In this case, our
   approach will fall squarely into Type B, as looping forever is not an option.

   \pause
   \vspace{\fill}

   So we need to accept being wrong sometimes. Let's see how.
 \end{frame}

 \sectionSlide{4}{Dataflow Analysis}

 \begin{frame}[fragile]
   \frametitle{Problems in Program Analysis}

   What kinds of questions are we interested in answering with program analysis?

   \pause
   \vspace{\fill}

   There are many, but a few that stand out are as follows:

   \pause
   \begin{itemize}
    \item is it possible for data from some source to reach some function? \pause
    \item is this code path unreachable? \pause
    \item is some computation redundant? \pause
    \item is an error state reachable?
   \end{itemize}
\end{frame}

 \begin{frame}[fragile]
   \frametitle{Problems in Program Analysis: Taint Analysis}

   The problem of \term{taint tracking} entails whether it is possible for data
   from some source to reach a specified function.

   \pause
   \vspace{\fill}

   \begin{pythoncodeblock}
      def foo():
        x = `input()`
        y = x.strip()
        /eval(y)/
   \end{pythoncodeblock}

   In this code example, we see that data from the call to \code{input()}
   goes through a \code{strip()}, which removes whitespace, and then reaches
   the call to the sensitive \code{eval()} function.

   \pause
   \vspace{\fill}

   In this case, we find that even though \code{input()} does not go
   \textit{directly} to \code{eval()}, it's still a security risk!
 \end{frame}

 \begin{frame}[fragile]
   \frametitle{Problems in Program Analysis: Dead Code Elimination}

   The problem of \term{dead code elimination} entails determining whether
   there is some code which cannot be reached under any path of execution.

   \pause
   \vspace{\fill}

   \begin{codeblock}[language=python]
      def foo():
        return 2

        `x = 4`
        `return x`
   \end{codeblock}

   \pause
   \vspace{\fill}

   In this example, we find that everything highlighted is completely unreachable,
   because there is a \code{return} statement above it that ends the function,
   before it ever gets to it. In this case, this is \term{dead code}, and
   can be eliminated.
 \end{frame}

 \begin{frame}[fragile]
   \frametitle{Problems in Program Analysis: Subexpression Elimination}

   The problem of \term{subexpression elimination} entails determining whether
   we are computing some expression twice, redundantly.

   \pause
   \vspace{\fill}

   \begin{codeblock}[language=python]
      def foo(n):
        y = `2 * n`
        z = y + `(2 * n)`
        return z
   \end{codeblock}

   \pause
   \vspace{\fill}

   In this example, we take in a number \code{n} and then compute \code{4 * n}
   in a really roundabout way. In particular, this is done by computing \code{2
   * n} twice!

   \pause
   \vspace{\fill}

   The code would be better rewritten as follows:

   \begin{codeblock}[language=python]
      def foo(n):
        y = 2 * n
        z = y + `y`
        return z
   \end{codeblock}
 \end{frame}

 \begin{frame}[fragile]
   \frametitle{Problems in Program Analysis: Code Reachability}

   The problem of \term{code reachability} has to do with determining whether or
   not there is a given series of inputs that can cause some code path to execute.

   \pause
   \vspace{\fill}

   \begin{codeblock}[language=python]
    def foo(person):
      if isTaylorSwiftFan(person):
        f()
      elif isLame(person):
        g()
      else:
        throw new Error("bad")
   \end{codeblock}

   \pause
   \vspace{\fill}

   In this case, we are interested in whether the \code{Error} can be reached.
   It's not clear, but a sophisticated program analysis tool could know that not
   being a Taylor Swift fan implies being a lame person, so the last condition is
   not actually reachable.
 \end{frame}

 \begin{frame}[fragile]
   \frametitle{Approximation}

   These problems all come with their own hardships. Due to what was discussed
   in the last section, they are all impossible to solve completely, but several
   of them are actually solvable via a Type B (allowed to be wrong) analysis,
   using the same method.

   \pause
   \vspace{\fill}

   A classic technique used in program analysis to obtain \textit{approximative}
   \footnote{Life hack: you can successfully replace "incorrect" with
   "approximative" in so many different places that it's hilarious.} answers
   in a finite amount of time is called \term{dataflow analysis}.

   \pause
   \vspace{\fill}

   Before I can define it to you, I must give you an analogy.

 \end{frame}

 \begin{frame}[fragile]
   \frametitle{An Analogy for Dataflow Analysis}

   Suppose you have a query you would like to solve on programs, which has many
   possible answers. Further suppose that the number of possible answers is finite.

   \pause
   \vspace{\fill}

   Suppose that you line them all up, one next to the other.

   \pause
   \vspace{\fill}

   \begin{center}
   \begin{tikzpicture}
    \node[box, node distance=1in] (A) {answer 1};
    \node[box, right of=A, node distance=1in] (B) {answer 2};
    \node[box, right of=B, node distance=1in] (C) {answer 3};
    \node[right of=C, node distance=1in] (D) {$\cdots$};
    \node[box, right of=D, node distance=1in] (E) {answer $n$};
   \end{tikzpicture}
    \end{center}

   \pause
   \vspace{\fill}

   Program analysis is hard because information can change a lot, infinitely much
   in fact, over the course of a program's run-time. For instance, suppose that
   the question is "how much memory is being used by the program?". You might pick an answer $n$,
   then move to answer $n - 2$, then move to a different answer $k$ altogether. It's
   possible to jump all around, in the limit of the program's execution.

   \pause
   \vspace{\fill}

   An observation can be made that, if you can order your answer in a way such that,
   over the course of your analysis, you only ever change your answer by moving
   right, you will always eventually terminate.
 \end{frame}

 \begin{frame}[fragile]
   \frametitle{Monotonicity}

   This is a roundabout way of describing what is known as a \term{monotonic function},
   which is a function which always "increases", according to some proper notion of
   "increases". In this case, our monotonic function also has an upper limit, i.e.
   a point beyond which it can no longer grow.

   \pause
   \vspace{\fill}

   For dataflow analysis, we will make use of this kind of analysis to iterate over
   our control-flow graph, constantly updating our answer, but only in a way that
   "increases", and eventually caps out. If we can do that, then we will
   guarantee that we will terminate.

   \pause
   \vspace{\fill}

   This also usually makes our answers sometimes wrong, though.
 \end{frame}

 \begin{frame}[fragile]
   \frametitle{Refresher: Control-flow Graphs}

   To understand dataflow analysis, first we need a really quick refresher on
   control-flow graphs.

   \pause
   \vspace{\fill}

   A \term{control-flow graph} is a kind of structure which basically describes all
   the paths that a program might take. When an instruction is \textit{always}
   followed by another, then they appear in a single "block" together. When the
   instructions might \textit{branch} (that is, go one place or another), then
   arrows point towards the blocks that it might lead to.

   \pause
   \vspace{\fill}

  \begin{center}
    \begin{minipage}{0.35\textwidth}
      \begin{pythoncodeblock}
        /x = True/
        /y = 2/

        /if x:/
          ?x = False?
        else:
          &y = 3&
      \end{pythoncodeblock}
    \end{minipage}
    \pause
    \begin{minipage}{0.6\textwidth}
      \begin{center}
      \begin{tikzpicture}
          [every node/.style={rectangle, draw}]
        \node[fill=bgBlue, text width=1.5in] (A) {
          \code{x = True} \strut \\
          \code{y = 2}\strut \\
          \code{if x:}\strut
        };
        \node[below of=A, node distance=1in, xshift=0.2in, fill=bgOrange] (B) {
          \code{x = False}\strut
        };
        \node[right of=B, node distance=1in, yshift=0.3in, fill=red!45] (C) {
          \code{y = 3}\strut
        };
        \draw[-stealth] (A) to [out=270, in=160] (B.west);
        \draw[-stealth] (A) to [out=270, in=180] (C);
      \end{tikzpicture}
      \end{center}
    \end{minipage}
  \end{center}
 \end{frame}

 \begin{frame}[fragile]
   \frametitle{Refresher: Control-flow Graphs}

   Note that control-flow graphs can get pretty massive and complicated.

   \vspace{\fill}

   The important takeaway here is that \textbf{control-flow graphs only predict
   program behavior}, they do not dictate it. Because the arrows just show where
   a program \textit{might} go, it doesn't actually tell us anything that the
   original program didn't.

   \vspace{\fill}

   The second important takeaway is that \textbf{there is always a way to turn a
   given program to a control-flow graph}. Let's not worry about how that happens.
 \end{frame}

 \begin{frame}[fragile]
   \frametitle{A Dataflow Example}

   \begin{center}
     \begin{minipage}{0.65\textwidth}
       \raggedright
       For instance, consider the following control-flow graph\footnotemark[13]:

       \vspace{10pt}

       We would like to perform an analysis known as \term{constant propagation}
       on it, by noting which variables are \term{constant}, i.e. never changed over
       the duration of the program.

       \vspace{10pt}

       How do we do this? We simply march forward through the CFG, and noting
       down which variables are constant as we see them, starting with the
       empty set.
     \end{minipage}
     \begin{minipage}{0.3\textwidth}
       \centering
       \begin{tikzpicture}[node distance=0.8in]
         \node[box, text width=1in] (N1) {
           \code{x = 1}\strut\\
           \code{y = 2}\strut
         };

         \node[box, below of=N1] (N2) {
           \code{x = x + 2}
         };

         \node[box, below of=N2] (N3) {
           \code{return x}
         };

         \draw[-stealth] (N1) -- (N2);
         \draw[-stealth] (N2.south west) to [out = 210, in=150, looseness=2] (N2.north west);
         \draw[-stealth] (N2) -- (N3);
       \end{tikzpicture}
     \end{minipage}
   \end{center}

   \footnotetext[13]{
Note that this
       one has a self-loop. This can happen, for example, from a \code{while} loop.
   }

 \end{frame}

 \begin{frame}[fragile]
   \frametitle{A Dataflow Example}

   \begin{center}
     \begin{minipage}{0.55\textwidth}
       \raggedright
       So for instance, first we traverse the entering block,
       by simply penciling in \code{x} and \code{y} as we see
       them get assigned to constants.

       \vspace{10pt}

       Once we finish, we now have the \term{out-set} for the
       first block, which we can then use to determine the other
       blocks.
     \end{minipage}
     \hfill
     \begin{minipage}{0.4\textwidth}
       \centering
       \begin{tikzpicture}[node distance=0.8in]
         \node[box, text width=1in, hlbg] (N1) {
           \code{x = 1}\strut\\
           \code{y = 2}\strut
         };

         \node[above of=N1, node distance=0.5cm, xshift=1.2in] {$\{\}$};
         \node[xshift=1.2in] at (N1) {$\{\code{x} \mapsto \code{1}\}$};
         \node[below of=N1, node distance=0.5cm, xshift=1.2in] {$\{\code{x} \mapsto \code{1}, \code{y} \mapsto \code{2}\}$};

         \node[box, below of=N1] (N2) {
           \code{x = x + 2}
         };

         \node[box, below of=N2] (N3) {
           \code{return x}
         };

         \draw[-stealth] (N1) -- (N2);
         \draw[-stealth] (N2.south west) to [out = 210, in=150, looseness=2] (N2.north west);
         \draw[-stealth] (N2) -- (N3);
       \end{tikzpicture}
     \end{minipage}
   \end{center}

 \end{frame}

 \begin{frame}[fragile]
   \frametitle{A Dataflow Example}

   \begin{center}
     \begin{minipage}{0.55\textwidth}
       \raggedright
       After following the highlighted edge, we end up at the second block. Since
       we have some information about what variables are constant, we can carry that
       information here.

       \vspace{10pt}

       Then, we see that \code{x} is incremented by two, and thus must be constant
       at \code{3} at the conclusion of the block.

     \end{minipage}
     \hfill
     \begin{minipage}{0.4\textwidth}
       \centering
       \begin{tikzpicture}[node distance=0.8in]
         \node[box, text width=1in] (N1) {
           \code{x = 1}\strut\\
           \code{y = 2}\strut
         };

         \node[above of=N1, node distance=0.5cm, xshift=1.2in] {$\{\}$};
         \node[xshift=1.2in] at (N1) {$\{\code{x} \mapsto \code{1}\}$};
         \node[below of=N1, node distance=0.5cm, xshift=1.2in] {$\{\code{x} \mapsto \code{1}, \code{y} \mapsto \code{2}\}$};

         \node[box, below of=N1, hlbg] (N2) {
           \code{x = x + 2}
         };

         \node[above of=N2, node distance=0.5cm, xshift=1.2in] {$\{\code{x} \mapsto \code{1}, \code{y} \mapsto \code{2}\}$};
         \node[below of=N2, node distance=0.5cm, xshift=1.2in] {$\{\code{x} \mapsto \code{3}, \code{y} \mapsto \code{2}\}$};

         \node[box, below of=N2] (N3) {
           \code{return x}
         };

         \draw[-stealth, hledge] (N1) -- (N2);
         \draw[-stealth] (N2.south west) to [out = 210, in=150, looseness=2] (N2.north west);
         \draw[-stealth] (N2) -- (N3);
       \end{tikzpicture}
     \end{minipage}
   \end{center}

 \end{frame}

 \begin{frame}[fragile]
   \frametitle{A Dataflow Example}

   \begin{center}
     \begin{minipage}{0.55\textwidth}
       \raggedright

       But, now we need to follow the self-loop! Something weird happens here.

       \vspace{10pt}

       There are two conflicting out-sets that are going in to the second block.
       One is the one we just computed,
       $\{\code{x} \mapsto \code{3}, \code{y} \mapsto \code{1}\}$, from the
       output of the second block itself. The other is
       $\{\code{x} \mapsto \code{1}, \code{y} \mapsto \code{1}\}$, from the
       original out-set from the first block.

       \vspace{10pt}

       This means we have a conflict. \code{x} is constant, but at two
       different values, coming in to the second block.
     \end{minipage}
     \hfill
     \begin{minipage}{0.4\textwidth}
       \centering
       \begin{tikzpicture}[node distance=0.8in]
         \node[box, text width=1in] (N1) {
           \code{x = 1}\strut\\
           \code{y = 2}\strut
         };

         \node[above of=N1, node distance=0.5cm, xshift=1.2in] {$\{\}$};
         \node[xshift=1.2in] at (N1) {$\{\code{x} \mapsto \code{1}\}$};
         \node[below of=N1, node distance=0.5cm, xshift=1.2in] {$\{{\color{hlset}\code{x} \mapsto \code{1}}, \code{y} \mapsto \code{2}\}$};

         \node[box, below of=N1, hlbg] (N2) {
           \code{x = x + 2}
         };

         \node[above of=N2, node distance=0.5cm, xshift=1.2in] {$\{\code{x} \mapsto \code{1}, \code{y} \mapsto \code{2}\}$};
         \node[below of=N2, node distance=0.5cm, xshift=1.2in] {$\{{\color{hlset}\code{x} \mapsto \code{3}}, \code{y} \mapsto \code{2}\}$};

         \node[box, below of=N2] (N3) {
           \code{return x}
         };

         \draw[-stealth] (N1) -- (N2);
         \draw[-stealth, hledge] (N2.south west) to [out = 210, in=150, looseness=2] (N2.north west);
         \draw[-stealth] (N2) -- (N3);
       \end{tikzpicture}
     \end{minipage}
   \end{center}
 \end{frame}

 \begin{frame}[fragile]
   \frametitle{A Dataflow Example}

   \begin{center}
     \begin{minipage}{0.55\textwidth}
       \raggedright

       This must mean that \code{x} is not constant after all.

       \vspace{10pt}

       So we set the value of \code{x} to instead be $\top$, which means
       "not constant". Note that this is different than \code{x} not having
       a value, which denotes not knowing if it is constant or not.
     \end{minipage}
     \hfill
     \begin{minipage}{0.4\textwidth}
       \centering
       \begin{tikzpicture}[node distance=0.8in]
         \node[box, text width=1in] (N1) {
           \code{x = 1}\strut\\
           \code{y = 2}\strut
         };

         \node[above of=N1, node distance=0.5cm, xshift=1.2in] {$\{\}$};
         \node[xshift=1.2in] at (N1) {$\{\code{x} \mapsto \code{1}\}$};
         \node[below of=N1, node distance=0.5cm, xshift=1.2in] {$\{\code{x} \mapsto \code{1}, \code{y} \mapsto \code{2}\}$};

         \node[box, below of=N1, hlbg] (N2) {
           \code{x = x + 2}
         };

         \node[above of=N2, node distance=0.5cm, xshift=1.2in] {$\{\code{x} \mapsto {\color{hlset}\top}, \code{y} \mapsto \code{2}\}$};
         \node[below of=N2, node distance=0.5cm, xshift=1.2in] {$\{\code{x} \mapsto {\color{hlset}\top}, \code{y} \mapsto \code{2}\}$};

         \node[box, below of=N2] (N3) {
           \code{return x}
         };

         \draw[-stealth] (N1) -- (N2);
         \draw[-stealth] (N2.south west) to [out = 210, in=150, looseness=2] (N2.north west);
         \draw[-stealth] (N2) -- (N3);
       \end{tikzpicture}
     \end{minipage}
   \end{center}
 \end{frame}

 \begin{frame}[fragile]
   \frametitle{A Dataflow Example}

   \begin{center}
     \begin{minipage}{0.55\textwidth}
       \raggedright

       Then, once we are assured that everything looks good, we can proceed
       to the final block, where we observe that we return \code{x} at a
       non-constant value.

       \vspace{10pt}

       This means that we cannot optimize the return value of this function
       after all.

       \vspace{10pt}

       Different story if we had returned \code{y}!
     \end{minipage}
     \hfill
     \begin{minipage}{0.4\textwidth}
       \centering
       \begin{tikzpicture}[node distance=0.8in]
         \node[box, text width=1in] (N1) {
           \code{x = 1}\strut\\
           \code{y = 2}\strut
         };

         \node[above of=N1, node distance=0.5cm, xshift=1.2in] {$\{\}$};
         \node[xshift=1.2in] at (N1) {$\{\code{x} \mapsto \code{1}\}$};
         \node[below of=N1, node distance=0.5cm, xshift=1.2in] {$\{\code{x} \mapsto \code{1}, \code{y} \mapsto \code{2}\}$};

         \node[box, below of=N1] (N2) {
           \code{x = x + 2}
         };

         \node[above of=N2, node distance=0.5cm, xshift=1.2in] {$\{\code{x} \mapsto \top, \code{y} \mapsto \code{2}\}$};
         \node[below of=N2, node distance=0.5cm, xshift=1.2in] {$\{\code{x} \mapsto \top, \code{y} \mapsto \code{2}\}$};

         \node[box, below of=N2, hlbg] (N3) {
           \code{return x}
         };

         \node[above of=N3, node distance=0.5cm, xshift=1.1in] {$\{\code{x} \mapsto \top, \code{y} \mapsto \code{2}\}$};
         \node[below of=N3, node distance=0.5cm, xshift=1.1in] {$\{\code{x} \mapsto \top, \code{y} \mapsto \code{2}\}$};

         \draw[-stealth] (N1) -- (N2);
         \draw[-stealth] (N2.south west) to [out = 210, in=150, looseness=2] (N2.north west);
         \draw[-stealth] (N2) -- (N3);
       \end{tikzpicture}
     \end{minipage}
   \end{center}
 \end{frame}

 \begin{frame}[fragile]
   \frametitle{Generality of Dataflow}

   This is a contrived example, but the really interesting thing is that
   dataflow analysis works for \textit{any} control-flow graph.

   \pause
   \vspace{\fill}

   The process seemed somewhat silly, as we could determine with our eyes that
   \code{x} was non-constant, but for very complicated control-flow graphs this
   is not an obvious fact at all. Programmatically, we can still run this same
   analysis, however.

   \pause
   \vspace{\fill}

   This analysis is also guaranteed to terminate, due to the monotonic reasons
   we stated earlier. The actual reason for this is that dataflow analysis
   strictly traverses up a \term{lattice}.
 \end{frame}

 \begin{frame}[fragile]
   \frametitle{Constant Definitions Lattice}

   \begin{center}
     \begin{tikzpicture}
       \node[bhex, fill=bgOrange] (T) {$\{\code{x} \mapsto \top, \code{y} \mapsto \top \}$};
       \node[bhex, fill=bgRed, below of=T, node distance=2.5in] (B) {$\{\}$};

       \node[bhex, below of=T, node distance=0.625in, xshift=-1in] (U1)
         {$\{\code{x} \mapsto n, \code{y} \mapsto \top\}$};
       \node[bhex, fill=bgPurple, below of=T, node distance=1.25in] (U2)
         {$\{\code{x} \mapsto n, \code{y} \mapsto n'\}$};
       \node[bhex, below of=T, node distance=0.625in, xshift=1in] (U3)
         {$\{\code{x} \mapsto \top, \code{y} \mapsto n' \}$};

       \node[bhex, above of=B, node distance=1in, xshift=-1.5in] (L1)
         {$\{\code{y} \mapsto \top\}$};
       \node[bhex, above of=B, node distance=0.625in, xshift=-0.5in] (L2)
         {$\{\code{x} \mapsto n\}$};
       \node[bhex, above of=B, node distance=0.625in, xshift=0.5in] (L3)
         {$\{\code{y} \mapsto n'\}$};
       \node[bhex, above of=B, node distance=1in, xshift=1.5in] (L4)
         {$\{\code{x} \mapsto \top \}$};

       \draw[-stealth] (B) -- (L1);
       \draw[-stealth] (B) -- (L2);
       \draw[-stealth] (B) -- (L3);
       \draw[-stealth] (B) -- (L4);

       \draw[-stealth] (L1) -- (U1);

       \draw[-stealth] (L4) -- (U3);

       \draw[-stealth] (L2) -- (U2);

       \draw[-stealth] (L3) -- (U2);

       \draw[-stealth] (U1) -- (T);
       \draw[-stealth] (U3) -- (T);

       \draw[-stealth] (U2) -- (U1);
       \draw[-stealth] (U2) -- (U3);

       \draw[-stealth] (U2) -- (U1);
       \draw[-stealth] (U2) -- (U3);

       \draw[-stealth] (L3) to [out=150, in=-10] (L1);
       \draw[-stealth] (L2) to [out=30, in=190] (L4);

     \end{tikzpicture}
   \end{center}
 \end{frame}


 \begin{frame}[fragile]
   \frametitle{Climbing the Lattice}

   The diagram looks scary, but the key thing is just that it
   assigns each variable a value of either no value, any constant $n$ or
   $n'$, or $\top$, which means "not constant".

   \pause
   \vspace{\fill}

   \begin{center}
   \begin{tikzpicture}
    \node[box, node distance=1.8in, text width=1.5in] (A) {assigned 0 times};
    \node[box, right of=A, node distance=1.8in, text width=1.5in] (B) {assigned once};
    \node[box, right of=B, node distance=1.8in, text width=1.5in] (C) {assigned more than once ($\top$)};

    \draw[-stealth] (A) -- (B);
    \draw[-stealth] (B) -- (C);
   \end{tikzpicture}
    \end{center}



   \pause
   \vspace{\fill}

   Edges go from sets to ones which have either added a new variable at
   a constant, or that have upgraded a variable from a constant to the
   not-a-constant symbol $\top$. This represents the gaining of information,
   of either a variable being declared as a constant, or a variable being
   discovered as not-a-constant.
   \pause
   \begin{itemize}
     \item
       For instance, we started at $\{\}$, and then
       went to $\{ \code{x} \mapsto \code{1} \}$ when we read \code{x = 1}.
       \pause
     \item
       Or, we went from
         $\{\code{x} \mapsto \code{1}, \code{y} \mapsto \code{2}\}$
       to
         $\{\code{x} \mapsto \top, \code{y} \mapsto \code{2}\}$
       upon seeing that \code{x} was set as two different constants, along
       two different paths to the block.
   \end{itemize}
 \end{frame}



 \begin{frame}[fragile]
   \frametitle{Monotonicity in Constant Definitions}

   The main thing to take away here is that \textit{every arrow goes up}.
   We talked earlier about putting answers on a number line, which is
   represented in terms of the lattice's height, here. No matter what
   edge you pick, you go up some amount, which means we must terminate.

   \pause
   \vspace{\fill}

   We can't go down, because it's impossible to update your worldview to
   either remove a variable from the set, or to set a variable from not-constant
   to constant. Once you know something about a variable, you can't go back.
 \end{frame}

 \begin{frame}[fragile]
   \frametitle{Approximation of Dataflow}
   \begin{center}
     \begin{minipage}{0.65\textwidth}
       \raggedright
       I mentioned earlier that this analysis is necessarily wrong, however.

       \vspace{10pt}

       The reason comes out of the fact that the control-flow graph is just how
       the execution of the program \textit{might} go. In reality, it's quite
       possible that at runtime, we never enter the self-loop, meaning that
       \code{x} really is constant.

       \vspace{10pt}

       But, without running the program, we have no way of knowing, so we
       assume that \code{x} \textit{is} updated at some point. This makes
       our knowledge necessarily possibly wrong, but a good approximation.
     \end{minipage}
     \begin{minipage}{0.3\textwidth}
       \centering
       \begin{tikzpicture}[node distance=0.8in]
         \node[box, text width=1in] (N1) {
           \code{x = 1}\strut\\
           \code{y = 2}\strut
         };

         \node[box, below of=N1] (N2) {
           \code{x = x + 2}
         };

         \node[box, below of=N2] (N3) {
           \code{return x}
         };

         \draw[-stealth] (N1) -- (N2);
         \draw[-stealth] (N2.south west) to [out = 210, in=150, looseness=2] (N2.north west);
         \draw[-stealth] (N2) -- (N3);
       \end{tikzpicture}
     \end{minipage}
   \end{center}
 \end{frame}

 \begin{frame}[fragile]
   \frametitle{On Dataflow Analysis}

   This was a really simple example that I hoped you might be able to understand.

   \pause
   \vspace{\fill}

   Dataflow analysis in general is a very powerful technique, however, and admits
   many other analyses, many of which are quite useful. These include:
   \pause
   \begin{itemize}
     \item \term{available expressions} - is there a definition of this exact
     expression already at this program point? useful in optimizing away
     redundant computations. \pause
     \item \term{reaching definitions} - what definitions of a variable can
     reach a given program point? useful in building use-def chains (i.e.
     "goto definition" in IDEs) \pause
     \item \term{liveness analysis} - what variables might actually be needed
     in the future? useful in eliminating dead code. \pause
     \item \term{taint tracking} - can data from undesirable sources reach
     some sensitive program point? \textbf{very} useful in security applications.
   \end{itemize}

   \pause
   \vspace{\fill}

   The rabbit hole goes deep. This one is simple, but this is a bread-and-butter
   technique in program analysis.
 \end{frame}

 \begin{frame}[fragile]
   \frametitle{Dataflow Analysis and Semantics}

   Dataflow analysis is a foundational technique that forms the bread and butter
   of the program analysis toolbox.

   \pause
   \vspace{\fill}

   In a sense, though, it can be viewed as a "lower level" analysis. We elided a lot
   of details in this presentation, but a proper dataflow analysis often entails breaking
   down a program into its most base representation. This tends to lose a lot of the
   original structure of the code -- nobody writes a CFG when programming.

   \pause
   \vspace{\fill}

   What if we could somehow stay at the level of the programmer when trying to do our
   program analysis? What if we could preserve the semantics and structure of the
   program while we look at it?

   \pause
   \vspace{\fill}

   It turns out, this will be a pretty effective approach to program analysis.

   \pause
   \vspace{\fill}

   This is the story of Semgrep.
 \end{frame}

 \sectionSlide{5}{Semgrep}

\begin{comment}
 \begin{frame}[fragile]
   \frametitle{Flavors of Program Analysis}

   It is difficult to give an accurate account of the state of cybersecurity within a
   single lecture (while talking about other things), but the main important thing to
   realize is that security is often the least concern of many developers.

   Developers want to write code, and are judged upon, promoted due to, and attached
   to the code that they write. This only makes the ultimate goal, the security of
   software throughout all sectors, a hard one to accomplish.

   There are engineers specifically hired to ensure the security of certain company's
   code, there are departments dedicated towards that same practice. The cost of a
   security breach is high, and nobody wants to be the next Yahoo.

   What does that mean? It means that we are on an uphill slope of a slope which
   is already upside-down.
 \end{frame}

 \begin{frame}[fragile]
   \frametitle{Flavors of Program Analysis}

   In the cybersecurity space, there are two big important aspects to any
   code-scanning security tool (or \term{SAST tool}).

   Firstly, it should terminate within a reasonable amount of time. Code is
   usually scanned on each PR\footnote<2->{Meaning \term{pull request}, or a
   potential contribution to a software project.}, and needs to fit within a
   developer's normal workflow, meaning a longer scan time than a few minutes is
   completely unreasonable.

   Secondly, it should not be too noisy. Security tools are going to produce
   matches that are wrong -- either false positives or false negatives. That's
   a part of the job. But they can't be wrong \textit{too} often, or that
   also marks an undesirable tool.
 \end{frame}

 \begin{frame}[fragile]
   \frametitle{Flavors of Program Analysis}

   To that end, security becomes an interesting game where the goal of a SAST
   tool is to work, but almost invisibly. Nobody likes a noisy tool. When a
   SAST tool pipes up, it had better be for a good reason.

   \defBox{}{We call a reported error or region of interest from a code-scanning
   tool a \term{finding}.}

   But another component is that a SAST tool should try to work
   \textit{with the developer}. It should be clear why a finding occurred,
   how it can be fixed, and it should get in the developer's way as little
   as possible. That's the game of security -- be functional, but frictionless.
 \end{frame}
\end{comment}

 \begin{frame}[fragile]
   \frametitle{Refresher: Abstract Syntax Trees}

   Programs are text files. To seasoned practitioners, however, programs are trees
   \footnote{To a viewer who is not familiar, there is not really much I can do to
   catch someone up to speed on trees in minutes. Basically, they are a way of storing
   data when the data might have multiple "descendants".}
   \footnote{It's kind of similar to how light is both a particle and a wave.}.

   \pause
   \vspace{\fill}

   \begin{minipage}{0.65\textwidth}
    For instance, consider simple arithmetic expressions, which are a subset of
    programs. We might represent the arithmetic expression \code{(1 - 2) + 3} as:

    When we write an expression like \code{(1 - 2) + 3} out on paper, we really
    mean this tree-like thing. Our notation on paper is just a syntax to convey
    the meaning, which we call an \term{abstract syntax tree}.
   \end{minipage}
   \hspace{\fill}
   \begin{minipage}{0.3\textwidth}
    \begin{center}
      \begin{tikzpicture}
        \node[circ] {\code{+}}
          child{node[circ] {\code{-}}
            child{node[hex] {\code{1}}}
            child{node[hex] {\code{2}}}
          }
          child{node[hex] {\code{3}}}
        ;

      \end{tikzpicture}
    \end{center}
  \end{minipage}

\end{frame}

\begin{frame}[fragile]
  \frametitle{Refresher: Abstract Syntax Trees}

  \begin{minipage}{0.35\textwidth}
    \raggedright
    For programs, this gets even more hardcore, actually. Consider the program:
    \begin{pythoncodeblock}
      if x == 3:
        while(y):
          pass
      else:
        return 2
    \end{pythoncodeblock}
  \end{minipage}
  \pause
  \hspace{\fill}
  \begin{minipage}{0.6\textwidth}
    \begin{center}
      \begin{tikzpicture}
        \node[box] {\code{if}}
          child{node[circ, xshift=-0.5in] {\code{==}}
            child{node[hex, fill=bgBlue] {\code{x}}}
            child{node[hex] {\code{3}}}
          }
          child{node[box] {\code{while}}
            child{node[hex, fill=bgBlue] {\code{y}}}
            child{node[box] {\code{pass}\strut}}
          }
          child{node[box, xshift=0.5in] {\code{return}}
            child{node[hex] {\code{2}}}
          }
        ;

      \end{tikzpicture}
    \end{center}
  \end{minipage}

  \pause
  \vspace{\fill}

  We see that, if you squint at the left and right, you can see that the indentation
  structure of the left program is actually reflected in the right tree.
  For instance, the indented \code{while} and \code{return} appear as children to the
  larger \code{if} node.
\end{frame}


 \begin{frame}[fragile]
   \frametitle{Analysis at the ASTs}

   Before we delve deeper, it's worth developing: why are we raving on about ASTs?
   What's the point?

   \pause
   \vspace{\fill}

   Semgrep as a technology makes heavy use of this idea of ASTs, by performing a
   great deal of its analysis at the level of the tree, as opposed to going
   further "down", to the level of the CFG. Our goal is to understand, at a high
   level, what Semgrep is doing.

  \pause
  \vspace{\fill}

  In this next section, we hope to go over some diagrams that can let you understand,
  at the level of a picture, what the Semgrep engine does when it produces a finding.
\end{frame}

 \begin{frame}[fragile]
   \frametitle{Security and Software Engineering}

  Before we proceed, it's also worth noting some things about security and software engineering.

  \pause
  \vspace{\fill}

  \textbf{These things do not generally go together.}

  \pause
  \vspace{\fill}

  Software engineers are paid to write code, are judged on their code, and spend
  large amounts of time thinking about their code. What they are not necessarily,
  however, is educated in the ways of safety. So this tends to take a back seat.
 \end{frame}

 \begin{frame}[fragile]
   \frametitle{Being Wrong}

   For most software, being wrong is a hellish scenario that is too horrifying
   to even contemplate.

   \pause
   \vspace{\fill}

   For a program analysis tool, being wrong is Tuesday.

   \pause
   \vspace{\fill}

   Still, to the developers who are receiving security notifications on their
   pull requests, as well as the security engineers that are constantly hovering
   over their shoulders\footnote{This presentation was made on macOS Monterey, cry about it Jonathan.\footnotemark[16]}, it matters \textbf{how wrong} the tool
   usually is.

   \pause
   \vspace{\fill}

   The easier way to make change is not to preach at someone the right
   way to do things, but to make them want to do things the right way. A
   SAST tool needs to be fruitful, but frictionless.

   \pause
   \vspace{\fill}

   Why am I mentioning this?

   \footnotetext[16]{That's a joke.}
 \end{frame}

 \begin{frame}[fragile]
   \frametitle{Deduplicating Efforts}

   We want to be able to provide as little friction as possible, while producing
   the best possible results. In the case of Semgrep, \textbf{staying primarily
   at the language level provides numerous benefits}.

   \pause
   \vspace{\fill}

   For most program-analyzing tools, they will have their own specific ASTs for
   each language. A Python linter, for instance, will have a Python AST, a C compiler
   will have a C AST, so on and so forth.

   \pause
   \vspace{\fill}

   The usual workflow for this looks something like this:
   \begin{center}
   \begin{tikzpicture}
    \node[box, node distance=2in, text width=1.1in] (A) {source text};
    \node[box, right of=A, node distance=2in, text width=1.1in, fill=bgPurple] (B) {language-specific AST};
    \node[box, right of=B, node distance=2in, text width=1.1in, fill=bgBlue] (C) {findings, data, etc};

    \draw[-stealth] (A) -- (B) node[midway, above]{parsing} ;
    \draw[-stealth] (B) -- (C) node[midway, above]{analysis };
   \end{tikzpicture}
    \end{center}
\end{frame}

 \begin{frame}[fragile]
   \frametitle{A Pipeline for Program Analysis}

   So, for multiple language-specific program analysis tools, the picture looks like this:

   \pause
   \vspace{\fill}

   \begin{center}
     \begin{tikzpicture}
       \node[box] (N1) {Java code\strut};
       \node[box, right of=N1, node distance=1.5in] (N2) {Python code\strut};
       \node[box, right of=N2, node distance=1.5in] (N3) {C code\strut};

       \node[box, below of=N1, node distance=0.6in, fill=bgPurple] (M1) {Java AST};
       \node[box, below of=N2, node distance=0.6in, fill=bgPurple] (M2) {Python AST\strut};
       \node[box, below of=N3, node distance=0.6in, fill=bgPurple] (M3) {C AST};

       \node[box, below of=M1, node distance=0.6in, fill=bgBlue] (L1) {Java analysis\strut};
       \node[box, below of=M2, node distance=0.6in, fill=bgBlue] (L2) {Python analysis\strut};
       \node[box, below of=M3, node distance=0.6in, fill=bgBlue] (L3) {C analysis\strut};

       \draw[-stealth] (N1) -- (M1);
       \draw[-stealth] (N2) -- (M2);
       \draw[-stealth] (N3) -- (M3);

       \draw[-stealth] (M1) -- (L1);
       \draw[-stealth] (M2) -- (L2);
       \draw[-stealth] (M3) -- (L3);
     \end{tikzpicture}
   \end{center}

   \pause
   \vspace{\fill}

   This leads to a massive amount of duplicated efforts.
 \end{frame}

 \begin{frame}[fragile]
   \frametitle{The Essence of Programming Languages}

   \begin{minipage}{0.55\textwidth}
    \raggedright
   For as much as people complain about programming languages\footnotemark[17], after some time with them you begin to realize that most of them look
   quite similar.

   \pause
   \vspace{10pt}

   The most controversial opinion to come out of this talk that
   you didn't expect: \textbf{the distinctions between most programming
   languages are pretty shallow}.
   \end{minipage}
   \begin{minipage}{0.43\textwidth}
      \raggedright
      \includegraphics[scale=0.25]{samepicture.png}
   \end{minipage}

   \footnotetext[17]{As have I.}
 \end{frame}

 \begin{frame}[fragile]
   \frametitle{It's All Just Code}

   Concretely, many languages have for loops, and many languages have exception
   raising and handling, and almost every language has functions, ways to define
   variables, and ways to call functions.

   \pause
   \vspace{\fill}

   Now, every language has its own miniscule idiosyncrasies which set them apart,
   which makes writing something like a universal compiler a pipe dream. But what
   if we're working in an application where being wrong in small ways doesn't really
   matter?

   \pause
   \vspace{\fill}

   For program analysis tools, being wrong is Tuesday.
 \end{frame}

 \begin{frame}[fragile]
   \frametitle{A Pipeline for Semgrep Analysis}

   Semgrep is a code-scanning tool for over 24 language\footnote{Maybe over 22, once we
   kill C/C++ and Elixir.}.

   \pause
   \vspace{\fill}

   The way that it achieves this is that it turns every single language into the
   same data type, called the Generic AST, which is a smörgåsbord union of all
   the language features in pretty much every language.

   \pause
   \vspace{\fill}

   This means that its pipeline looks remarkably simpler:

   \begin{center}
     \begin{tikzpicture}
       \node[box] (N1) {Java code\strut};
       \node[box, right of=N1, node distance=1.5in] (N2) {Python code\strut};
       \node[box, right of=N2, node distance=1.5in] (N3) {C code\strut};

       \node[hex, below of=N2, node distance=0.6in, fill=bgRed] (M) {Generic AST};

       \node[box, below of=M, node distance=0.6in, fill=bgBlue] (L) {Analysis};

       \draw[-stealth] (N1) -- (M);
       \draw[-stealth] (N2) -- (M);
       \draw[-stealth] (N3) -- (M);

       \draw[-stealth] (M) -- (L);
     \end{tikzpicture}
   \end{center}
 \end{frame}

 \begin{frame}[fragile]
   \frametitle{Picking Battles}

   We see that staying at the AST level saves us work, which can be wrong at
   times due to imperfect translation, but ultimately doesn't matter for a
   program analysis tool. More importantly, it allows for cleaner code to
   be written, which translates to a better tool.

   \pause
   \vspace{\fill}

   When solving an impossible problem, you have to be pragmatic. Picking
   battles that you can win is essential.

   \pause
   \vspace{\fill}

   Semgrep takes a more "syntactic" approach to program analysis that gives
   benefits in terms of its development velocity as a technology. It also lends an
   advantage in its customizability and ease of use, which is the next thing to discuss.
 \end{frame}

 \begin{frame}[fragile]
   \frametitle{Refresher: Interpreted and Compiled Languages}

   It's worth also getting a picture of how Semgrep fits into the ecosystem of program analysis
   tools. Why is this choice, to only work at the level of source text, and require no further
   information, a competitive edge?

   \pause
   \vspace{\fill}

   The lifecycle of a program is varied. As a refresher, recall the notion of
   programming languages as \term{interpreted} or \term{compiled}.

   \pause
   \vspace{\fill}

   A program is a text file written in some kind of programming language. Some languages
   are written in \term{interpreted languages}, meaning that they can be executed (or \textit{run})
   immediately.

   \pause
   \vspace{\fill}

   Other languages are \term{compiled languages}, meaning that an executable version of
   the program must first be obtained by running a separate program called the
   \term{compiler}. This is also known as \term{compiling}, or \term{building} the program.
 \end{frame}



 \begin{frame}[fragile]
   \frametitle{An Americentric Analogy}

   For an Americentric analogy, you can think of a program in an interpreted language as
   a recipe written in English, which the average American can follow just by reading.
   A program in a compiled language is written in French, which, for the average American,
   must first be translated to English before using.

   \pause
   \vspace{\fill}

   \begin{center}
     \begin{tikzpicture}
       \node[hex, minimum size=1cm] (A) {English recipe};
       \node[hex, right of=A, node distance=2in, minimum size=1cm, fill=red!30!white] (B) {food};

       \draw[-stealth] (A) -- (B) node[midway, above] {cooking};
     \end{tikzpicture}

     \vspace{\fill}

     \begin{tikzpicture}
       \node[box, minimum size=1cm, fill=orange!40!white] (A) {French recipe};
       \node[hex, minimum size=1cm, right of=A, node distance=2.5in] (B) {English recipe};
       \node[hex, right of=B, node distance=2.5in, minimum size=1cm, fill=red!30!white] (C) {food};

       \draw[-stealth] (A) -- (B) node[midway, above] {translation};
       \draw[-stealth] (B) -- (C) node[midway, above] {cooking};
     \end{tikzpicture}
   \end{center}
 \end{frame}

 \begin{frame}[fragile]
   \frametitle{Refresher: Interpreted and Compiled Languages}

   Interpreted and compiled languages work similarly. By analogy, the comparison should be
   clear.

   \pause
   \vspace{\fill}

   \begin{center}
     \begin{tikzpicture}
       \node[hex, minimum size=1cm] (A) {program};
       \node[hex, right of=A, node distance=2in, minimum size=1cm, fill=red!30!white] (B) {result};

       \draw[-stealth] (A) -- (B) node[midway, above] {execution};
     \end{tikzpicture}

     \vspace{\fill}

     \begin{tikzpicture}
       \node[box, minimum size=1cm, fill=orange!40!white] (A) {program};
       \node[hex, minimum size=1cm, right of=A, node distance=2.5in] (B) {compiled program};
       \node[hex, right of=B, node distance=2.5in, minimum size=1cm, fill=red!30!white] (C) {result};

       \draw[-stealth] (A) -- (B) node[midway, above] {compilation};
       \draw[-stealth] (B) -- (C) node[midway, above] {execution};
     \end{tikzpicture}
   \end{center}
 \end{frame}

 \begin{frame}[fragile]
   \frametitle{The Multiplicity of Compiled Languages}

   Let's put interpreted languages to the side, for the moment.

   \pause
   \vspace{\fill}

   The picture is a nice one, but significantly simplified from reality. In truth, it's
   not that there is a single arrow leading to a single compiled program -- there are
   \textit{many} such arrows, depending on your particular compiler, build environment,
   etc.

   \pause
   \vspace{\fill}

   So the picture looks more like:
   \pause
  \begin{center}
    \begin{tikzpicture}
      \node[box, minimum size=0.6cm, fill=orange!40!white] (A) {program};
      \node[hex, minimum size=0.6cm, right of=A, node distance=2.5in] (B) {compiled program B};
      \node[hex, minimum size=0.6cm, above of=B, node distance=0.4in] (B1) {compiled program A};
      \node[hex, minimum size=0.6cm, below of=B, node distance=0.4in] (B2) {compiled program C};
      \node[above of=B1, node distance=0.4in, xshift=-1in] (B3) {\textellipsis};
      \node[below of=B2, node distance=0.4in, xshift=-1in] (B4) {\textellipsis};
      \node[hex, right of=B, node distance=2.5in, minimum size=0.6cm, fill=red!30!white] (C) {result B};
      \node[hex, right of=B1, node distance=2.5in, minimum size=0.6cm, fill=red!30!white] (C1) {result A};
      \node[hex, right of=B2, node distance=2.5in, minimum size=0.6cm, fill=red!30!white] (C2) {result C};

      \draw[-stealth] (A) -- (B) node[midway, above] {compilation};
      \draw[-stealth, postaction={decorate, decoration={markings, mark=at position 0.5 with {\node[above, xshift=0.15in] {\small compilation};}}}] (A.north) to [out=30, in=180] (B1.west);
      \draw[-stealth, postaction={decorate, decoration={markings, mark=at position 0.5 with {\node[above, xshift=0.15in] {\small compilation};}}}] (A.south) to [out=330, in=180] (B2.west);
      \draw[-stealth, postaction={decorate, decoration={markings, mark=at position 0.5 with {\node[above, xshift=-0.2in] {\small compilation};}}}] (A.north) to [out=60, in=180] (B3.west);
      \draw[-stealth, postaction={decorate, decoration={markings, mark=at position 0.5 with {\node[below, xshift=-0.2in] {\small compilation};}}}] (A.south) to [out=300, in=180] (B4.west);
      \draw[-stealth] (B) -- (C) node[midway, above] {\small execution};
      \draw[-stealth] (B1) -- (C1) node[midway, above] {\small execution};
      \draw[-stealth] (B2) -- (C2) node[midway, above] {\small execution};
    \end{tikzpicture}
  \end{center}
 \end{frame}

 \begin{frame}[fragile]
   \frametitle{SAST 1.0 and the Program Lifecycle}

   For many of the SAST 1.0 tools, they operate in the stage of the pipeline that is \textit{during compilation}.

   \pause
   \vspace{\fill}

   \begin{center}
   \includegraphics[scale=0.2]{sast1.0.png}
   \end{center}

   \pause
   \vspace{\fill}

   This means that they exist across \textit{all} of the different arrows of
   different ways to compile, and the different configurations which might exist at that point.

   \pause
   \vspace{\fill}

   This is not just a maintenance burden on the part of the tool, but also a burden on the
   user of the tool, which needs to be able to "hook in" the tool to that user's specific
   build concerns. This can be messy, error-prone, and intensive.
 \end{frame}

 \begin{frame}[fragile]
   \frametitle{Semgrep and Shifting Left}

   Semgrep operates slightly differently, by operating \textit{purely} upon the source
   text of the program. No additional configuration is needed -- the only things that it takes in
   are the program text, and what it should find.

   \pause
   \vspace{\fill}

   \begin{center}
   \includegraphics[scale=0.2]{semgrep!.png}
   \end{center}

   \pause
   \vspace{\fill}

   This means that Semgrep succeeds as a tool which is \textit{farther left}, in that it operates
   at a stage of the program lifecycle that is \textit{chronologically before} that of other tools.

   \pause
   \vspace{\fill}

   This lack of need for configuration and knowledge of build systems makes Semgrep incredibly easy
   to just drop in, compared to other tools. This is one benefit that Semgrep gets from a more strictly
   AST-based approach.
 \end{frame}

 \begin{frame}[fragile]
   \frametitle{Source-Level Errors}

   Many program errors aren't purely due to what happens
   at the level of assembly. Many of them are apparent from the source code,
   because they were written by a developer who was writing source code.

   \pause
   \vspace{\fill}

   For instance, consider the following Python code, with a function that
   has a default argument \code{to} set to \code{[]}:
   \begin{pythoncodeblock}
     def append_to(element, to=[]):
       to.append(element)
       return to
   \end{pythoncodeblock}

   \vspace{\fill}
   \pause

   What do you think this does?
 \end{frame}

 \begin{frame}[fragile]
   \frametitle{Default Dead Arguments}

   \begin{pythoncodeblock}
     def append_to(element, to=[]):
       to.append(element)
       return to
   \end{pythoncodeblock}

   \vspace{\fill}

   It turns out that the first time you call \code{append_to(2)}, you will
   get \code{[2]}. That's cool.

   \pause
   \vspace{\fill}

   The second time you call \code{append_to(2)}, you will get \code{[2, 2]}.

   \pause
   \vspace{\fill}

   Uh oh.
 \end{frame}

 \begin{frame}[fragile]
   \frametitle{Errors at the Source}

   It turns out this is due to the fact that Python default arguments are
   instantiated at function definition time, and so if it is mutated, that
   mutation persists through each call.

   \pause
   \vspace{\fill}

   This is an error that would be very easy to catch in code review. Unfortunately,
   there are tens of millions of developers in the world, and even a 99\%
   review success rate has bad implications.

   \pause
   \vspace{\fill}

   So at the end of the day, we might not \textit{need} to descend further down
   than the AST. Let's see how we can fix such bugs while staying above ground.
 \end{frame}

 \begin{frame}[fragile]
   \frametitle{Forms of Program Analysis}

   It turns out that program analysis can take many forms. We might want to use
   it for security purposes, we might want to use it for style checking, and we
   might want to use it to automate code review for things like the error we just
   saw.

   \pause
   \vspace{\fill}

   A common theme in program analysis is simply finding certain things in your
   source code. This might be something like an unused variable, a type error,
   or even a function call which might lead to a security vulnerability. We
   call such regions of interest a \term{finding}.

   \pause
   \vspace{\fill}

   Semgrep is aimed at being customizable to all of these use cases. This means
   that while that while some tools find security vulnerabilities, and some
   tools find style errors, the answer to what Semgrep finds is \textit{whatever
   you want}.
 \end{frame}

 \begin{frame}[fragile]
   \frametitle{A Semantic Grep}

   {\color{blue}\href{https://semgrep.dev/}{Semgrep}} stands for \textit{semantic grep}.

   \pause
   \vspace{\fill}

   If you've used the tool \code{grep} before, you know that it is really just a
   no-nonsense precursor to what is now Ctrl-F in most browsers.\footnote{Such as Arc,
   the browser that I am currently using.}

   \pause
   \vspace{\fill}

   But, it's limited in some ways. Suppose we wanted to find all instances of
   using the function \code{print} to debug, before we put the code up for
   review. What do you think is going to happen when we search for every instance
   of the literal substring "print"?

   \pause
   \vspace{\fill}

   The answer: \textbf{We are going to be absolutely inundated with results.} We'll
   find the string \code{print} if it occurs in comments, if it occurs in literal
   strings, and even if it's part of a larger method name. That doesn't make it what
   we were looking for, though!

   \pause
   \vspace{\fill}

   In other words, \code{grep} doesn't understand the \textit{meaning} of what we're
   looking for, it's just looking for sequences of characters. Semgrep does.
 \end{frame}

 \begin{frame}[fragile]
   \frametitle{Tree Search and Semgrep}

   How does Semgrep understand the semantics of the program? Recall that
   the literal text in a program just serve as a proxy that indicates the underlying
   tree of the program, which contains its real meaning.

   \pause
   \vspace{\fill}

   Semgrep doesn't do text search, it does \textit{tree search}. This means that
   instead of searching for a sequence of characters within a given program, it
   searches for a \textit{matching subtree} within a program's AST.

   \pause
   \vspace{\fill}

   This means that a user might specify a particular \term{pattern AST}, which refers
   to the AST that they would like to search for, and then Semgrep will search the
   \term{target AST} to find a subtree which matches it.

   \pause
   \vspace{\fill}

   The key innovation of Semgrep is that, although it does AST matching, the user
   doesn't need to be able to construct one, or even know what it is. Pattern ASTs
   are derived from patterns written \textit{in the source language}, meaning
   that they look like code.
 \end{frame}

 \begin{frame}[fragile]
   \frametitle{Matching, Take 1}


   \begin{center}
     \begin{minipage}[t][0.1\textheight][t]{\textwidth}
       {\color{blue}\href{https://semgrep.dev/playground/s/Pe9Dq}{So, our query might look like:}}
     \end{minipage}
     \begin{minipage}[t][0.6\textheight][t]{\textwidth}
       \begin{minipage}{0.38\textwidth}
         \centering
         \begin{tikzpicture}
           \node[hex] (A) {\code[keywordstyle=\color{black}]{print}\strut};
         \end{tikzpicture}
       \end{minipage}
       \begin{minipage}{0.58\textwidth}
         \centering
         \begin{tikzpicture}
           \node[box] (N1) {\code{Assign}\strut}
             child{node[circ, fill=bgRed, xshift=-0.3in] (N2) {\code{x}}}
             child{node[bbox, xshift=0.3in] (N3) {\code{Call}}
               child{node[hex, xshift=-1cm] (N4) {\code[keywordstyle=\color{black}]{print}\strut}}
               child{node[bbox, xshift=1cm] (N5) {\code{"let's print!"}\strut}}
             };
         \end{tikzpicture}
       \end{minipage}
     \end{minipage}
     \begin{minipage}[t][0.3\textheight][t]{\textwidth}
       \begin{minipage}{0.38\textwidth}
         \centering
         \begin{codeblock}
           print
         \end{codeblock}
       \end{minipage}
       \begin{minipage}{0.58\textwidth}
         \begin{codeblock}
           x = print("let's print!")
         \end{codeblock}
       \end{minipage}
     \end{minipage}
   \end{center}
 \end{frame}

 \begin{frame}[fragile]
   \frametitle{Matching, Take 1}

   \begin{center}
     \begin{minipage}[t][0.1\textheight][t]{\textwidth}
       {\color{blue}\href{https://semgrep.dev/playground/s/Pe9Dq}{So, our query might look like:}}
     \end{minipage}
     \begin{minipage}[t][0.6\textheight][t]{\textwidth}
       \begin{minipage}{0.38\textwidth}
         \centering
         \begin{tikzpicture}
           \node[hex, highlight] (A) {\code[keywordstyle=\color{black}]{print}\strut};
         \end{tikzpicture}
       \end{minipage}
       \begin{minipage}{0.58\textwidth}
         \centering
         \begin{tikzpicture}
           \node[box] (N1) {\code{Assign}\strut}
             child{node[circ, fill=bgRed, xshift=-0.3in] (N2) {\code{x}}}
             child{node[bbox, xshift=0.3in] (N3) {\code{Call}}
               child{node[hex, xshift=-1cm, highlight] (N4) {\code[keywordstyle=\color{black}]{print}\strut}}
               child{node[bbox, xshift=1cm] (N5) {\code{"let's print!"}\strut}}
             };
         \end{tikzpicture}
       \end{minipage}
     \end{minipage}
     \begin{minipage}[t][0.3\textheight][t]{\textwidth}
       \begin{minipage}{0.38\textwidth}
         \centering
         \begin{codeblock}
           print
         \end{codeblock}
       \end{minipage}
       \begin{minipage}{0.58\textwidth}
         \begin{codeblock}
           x = `print`("let's print!")
         \end{codeblock}
       \end{minipage}
     \end{minipage}
   \end{center}
 \end{frame}

 \begin{frame}[fragile]
   \frametitle{Syntax Aware}

   Notably, however, we avoid matching the instance of the literal string
   of "print" inside of the string that is being passed to its namesake function.

   \pause
   \vspace{\fill}

   This is because we are only matching at the granularity of nodes of the AST,
   not of the constituent text. Regular expressions, for instance, would be fooled
   by instances of "print" within comments or strings.

   \pause
   \vspace{\fill}

   When it comes to minimizing instances of program analysis tools being flat-out
   wrong, this is an invaluable help. Most program analysis tools are black-box
   applications that perform some magic to produce matches, without necessarily
   being understandable. The algorithm of Semgrep is really quite simple, and
   still manages to avoid false positives.
 \end{frame}


 \begin{frame}[fragile]
   \frametitle{Matching, Take 2}

   \begin{center}
     \begin{minipage}[t][0.21\textheight][t]{\textwidth}
       \raggedright
       {\color{blue}\href{https://semgrep.dev/playground/s/JDqJQ}{We can do a little more advanced of a query, too.}} Suppose that we were
       interested in finding all of the examples of a call to the function
       \code{ref}.

       \vspace{5pt}

       Well, we might write a pattern like \code{ref $\,\,\dollar$X}, which uses a
       \textit{metavariable} to bind to any sub-AST.
     \end{minipage}
     \begin{minipage}[t][0.515\textheight][t]{\textwidth}
       \begin{minipage}{0.48\textwidth}
         \centering
         \begin{tikzpicture}
           \node[bbox] (N1) {\code{Call}}
             child{node[bbox] (N2) {\code[keywordstyle=\color{black}]{foo}}}
             child{node[unknown] (N3) {\texttt{\$}\code{X}}}
           ;
         \end{tikzpicture}
       \end{minipage}
       \begin{minipage}{0.48\textwidth}
         \centering
         \begin{tikzpicture}
           \node[box] (N1) {\code{Assign}\strut}
             child{node[circ, fill=bgRed, xshift=-0.3in] (N2) {\code{y}}}
             child{node[bbox, xshift=0.3in] (N3) {\code{Call}}
               child{node[hex, xshift=-1cm] (N4) {\code[keywordstyle=\color{black}]{foo}}}
               child{node[bbox, xshift=1cm] (N5) {\code{2}}}
             };
         \end{tikzpicture}
       \end{minipage}
     \end{minipage}
     \begin{minipage}[t][0.3\textheight][t]{\textwidth}
       \begin{minipage}{0.48\textwidth}
         \centering
         \begin{codeblock}
           foo($\dollar$X)
         \end{codeblock}
       \end{minipage}
       \begin{minipage}{0.48\textwidth}
         \begin{codeblock}
           y = foo(2)
         \end{codeblock}
       \end{minipage}
     \end{minipage}
   \end{center}
 \end{frame}


 \begin{frame}[fragile]
   \frametitle{Matching, Take 2}

   \begin{center}
     \begin{minipage}[t][0.21\textheight][t]{\textwidth}
       \raggedright
       {\color{blue}\href{https://semgrep.dev/playground/s/JDqJQ}{Here, we bind
       the metavariable}} \tikz[baseline] \node[unknown, scale=0.7, anchor=base] {\texttt{\$}\code{X}}; to the node of \tikz[baseline] \node[bbox, anchor=base, scale=0.7]{2};, and then
       the entire pattern succeeds at matching the highlighted sub-tree of the
       target.
     \end{minipage}
     \begin{minipage}[t][0.515\textheight][t]{\textwidth}
       \begin{minipage}{0.48\textwidth}
         \centering
         \begin{tikzpicture}
           \node[bbox, highlight] (N1) {\code{Call}}
             child{node[bbox, highlight] (N2) {\code[keywordstyle=\color{black}]{foo}}}
             child{node[unknown, highlight] (N3) {\texttt{\$}\code{X}}}
           ;
         \end{tikzpicture}
       \end{minipage}
       \begin{minipage}{0.48\textwidth}
         \centering
         \begin{tikzpicture}
           \node[box] (N1) {\code{Assign}\strut}
             child{node[circ, fill=bgRed, xshift=-0.3in] (N2) {\code{y}}}
             child{node[bbox, xshift=0.3in, highlight] (N3) {\code{Call}}
               child{node[hex, xshift=-1cm, highlight] (N4) {\code[keywordstyle=\color{black}]{foo}}}
               child{node[bbox, xshift=1cm, highlight] (N5) {\code{2}}}
             };
         \end{tikzpicture}
       \end{minipage}
     \end{minipage}
     \begin{minipage}[t][0.3\textheight][t]{\textwidth}
       \begin{minipage}{0.48\textwidth}
         \centering
         \begin{codeblock}
           foo($\dollar$X)
         \end{codeblock}
       \end{minipage}
       \begin{minipage}{0.48\textwidth}
         \begin{codeblock}
           y = `foo(2)`
         \end{codeblock}
       \end{minipage}
     \end{minipage}
   \end{center}
 \end{frame}


 \begin{frame}[fragile]
   \frametitle{Matching, Take 3}

   \begin{center}
     \begin{minipage}[t][0.21\textheight][t]{\textwidth}
       \raggedright
       {\color{blue}\href{https://semgrep.dev/playground/s/5rW33}{We can also express more complicated queries.}} What if we want to find instances
       of a useless comparison, such as testing if an expression is equal to itself?
     \end{minipage}
     \begin{minipage}[t][0.515\textheight][t]{\textwidth}
       \begin{minipage}{0.48\textwidth}
         \centering
         \begin{tikzpicture}
           \node[circ] (N1) {\code{==}}
             child{node[unknown] (N2) {\code{$\dollar$X}}}
             child{node[unknown] (N3) {\code{$\dollar$X}}}
           ;
         \end{tikzpicture}
       \end{minipage}
       \begin{minipage}{0.48\textwidth}
         \centering
         \begin{tikzpicture}[level distance=1cm]
           \node[box] (N1) {\code{Assign}}
             child{node[hex, xshift=-0.3in] (N2) {\code{b}}}
             child{node[circ] (N3) {\code{==}}
               child{node[bbox, xshift=-0.35in] (N4) {\code{Call}}
                 child{node[hex] (N5) {\code{f}}}
                 child{node[bbox] (N6) {\code{2}}}
               }
               child{node[bbox, xshift=0.35in] (N7) {\code{Call}}
                 child{node[hex] (N8) {\code{f}}}
                 child{node[bbox] (N9) {\code{2}}}
               }
             }
           ;
         \end{tikzpicture}
       \end{minipage}
     \end{minipage}
     \begin{minipage}[t][0.3\textheight][t]{\textwidth}
       \begin{minipage}{0.48\textwidth}
         \centering
         \begin{codeblock}
           $\dollar$X == $\dollar$X
         \end{codeblock}
       \end{minipage}
       \begin{minipage}{0.48\textwidth}
         \begin{codeblock}
           b = (f(2) == f(2))
         \end{codeblock}
       \end{minipage}
     \end{minipage}
   \end{center}
 \end{frame}

 \begin{frame}[fragile]
   \frametitle{Matching, Take 3}

   \begin{center}
     \begin{minipage}[t][0.21\textheight][t]{\textwidth}
       \raggedright
       {\color{blue}\href{https://semgrep.dev/playground/s/GdR7n}{Here, because
       we reuse the metavariable}} \tikz[baseline] \node[unknown, anchor=base, scale=0.7] {\code{$\dollar$X}}; twice, the two instances
       \term{unify}, meaning that the ASTs that each binds to must be the same.
       In this case, that works.
     \end{minipage}
     \begin{minipage}[t][0.515\textheight][t]{\textwidth}
       \begin{minipage}{0.48\textwidth}
         \centering
         \begin{tikzpicture}
           \node[circ] (N1) {\code{==}}
             child{node[unknown] (N2) {\code{$\dollar$X}}}
             child{node[unknown] (N3) {\code{$\dollar$X}}}
           ;
         \end{tikzpicture}
       \end{minipage}
       \begin{minipage}{0.48\textwidth}
         \centering
         \begin{tikzpicture}[level distance=1cm]
           \node[box] (N1) {\code{Assign}}
             child{node[hex, xshift=-0.3in] (N2) {\code{b}}}
             child{node[circ, highlight] (N3) {\code{==}}
               child{node[bbox, xshift=-0.35in, highlight] (N4) {\code{Call}}
                 child{node[hex, highlight] (N5) {\code{f}}}
                 child{node[bbox, highlight] (N6) {\code{2}}}
               }
               child{node[bbox, xshift=0.35in, highlight] (N7) {\code{Call}}
                 child{node[hex, highlight] (N8) {\code{f}}}
                 child{node[bbox, highlight] (N9) {\code{2}}}
               }
             }
           ;
         \end{tikzpicture}
       \end{minipage}
     \end{minipage}
     \begin{minipage}[t][0.3\textheight][t]{\textwidth}
       \begin{minipage}{0.48\textwidth}
         \centering
         \begin{codeblock}
           $\dollar$X == $\dollar$X
         \end{codeblock}
       \end{minipage}
       \begin{minipage}{0.48\textwidth}
         \begin{codeblock}
           b = `(f(2) == f(2))`
         \end{codeblock}
       \end{minipage}
     \end{minipage}
   \end{center}
 \end{frame}

 \begin{frame}[fragile]
   \frametitle{Matching, Take 3}

   \begin{center}
     \begin{minipage}[t][0.21\textheight][t]{\textwidth}
       \raggedright
       {\color{blue}\href{https://semgrep.dev/playground/s/Rer0K}{But in this case}}, there would be no match, because the {\color{blue}blue} and
       {\color{magenta}magenta} subtrees do not match. In particular, \tikz[baseline] \node[hex, scale=0.7, anchor=base] {\code{f}};
       does not match \tikz[baseline] \node[hex, scale=0.7, anchor=base] {\code{g}};.

     \end{minipage}
     \begin{minipage}[t][0.515\textheight][t]{\textwidth}
       \begin{minipage}{0.48\textwidth}
         \centering
         \begin{tikzpicture}
           \node[circ] (N1) {\code{==}}
             child{node[unknown] (N2) {\code{$\dollar$X}}}
             child{node[unknown] (N3) {\code{$\dollar$X}}}
           ;
         \end{tikzpicture}
       \end{minipage}
       \begin{minipage}{0.48\textwidth}
         \centering
         \begin{tikzpicture}[level distance=1cm]
           \node[box] (N1) {\code{Assign}}
             child{node[hex, xshift=-0.3in] (N2) {\code{b}}}
             child{node[circ] (N3) {\code{==}}
               child{node[bbox, xshift=-0.35in, highlight2] (N4) {\code{Call}}
                 child{node[hex, highlight2] (N5) {\code{f}}}
                 child{node[bbox, highlight2] (N6) {\code{2}}}
               }
               child{node[bbox, xshift=0.35in, highlight3] (N7) {\code{Call}}
                 child{node[hex, highlight3] (N8) {\code{g}}}
                 child{node[bbox, highlight3] (N9) {\code{2}}}
               }
             }
           ;
         \end{tikzpicture}
       \end{minipage}
     \end{minipage}
     \begin{minipage}[t][0.3\textheight][t]{\textwidth}
       \begin{minipage}{0.48\textwidth}
         \centering
         \begin{codeblock}
           $\dollar$X == $\dollar$X
         \end{codeblock}
       \end{minipage}
       \begin{minipage}{0.48\textwidth}
         \begin{codeblock}
           b = (f(2) == g(2))
         \end{codeblock}
       \end{minipage}
     \end{minipage}
   \end{center}
 \end{frame}

 \begin{frame}[fragile]
   \frametitle{A Default Analysis}

   Let's apply this to a very concrete example -- the erroneous default list parameter we
   saw earlier, in Python! A pattern for this might look like:

   \pause
   \vspace{\fill}

   We can write a Semgrep query which checks for precisely that!

   \pause
   \vspace{\fill}

   \begin{minipage}{0.45\textwidth}
    \begin{pythoncodeblock}
      def $\dollar$F(..., $\dollar$ARG=[], ...):
        ...
    \end{pythoncodeblock}
   \end{minipage}
   \pause
    \begin{minipage}{0.53\textwidth}
      \centering
      \begin{tikzpicture}[level distance=1cm]
        \node[bbox] (N1) {\code{Def}}
          child{node[unknown, xshift=-0.25in] {\code{$\dollar$F}}}
          child{
            node[bbox, fill=bgRed, yshift=-0.1in] {\code{Args}\strut}
              child{node[unknown2, xshift=-0.4in] {\code{...}}}
              child{node[bbox, yshift=-0.1in] {\code{DefaultParam}}
                child{node[unknown] {\code{$\dollar$ARG}}}
                child{node[hex] {\code{[]}}}
              }
              child{node[unknown2, xshift=0.4in] {\code{...}}}
          }
          child{node[unknown2, xshift=0.25in] {\code{...}}}
        ;
      \end{tikzpicture}
    \end{minipage}
 \end{frame}

 \begin{frame}[fragile]
   \frametitle{A Default Analysis}

   \begin{center}
    %  \begin{minipage}[t][0.21\textheight][t]{\textwidth}
    %    \raggedright
    %    {\color{blue}\href{https://semgrep.dev/playground/s/Rer0K}{But in this case}}, there would be no match, because the {\color{blue}blue} and
    %    {\color{magenta}magenta} subtrees do not match. In particular, \tikz[baseline] \node[hex, scale=0.7, anchor=base] {\code{f}};
    %    does not match \tikz[baseline] \node[hex, scale=0.7, anchor=base] {\code{g}};.

    %  \end{minipage}
     \begin{minipage}[t][0.65\textheight][t]{\textwidth}
       \begin{minipage}{0.44\textwidth}
         \centering
          \begin{tikzpicture}[level distance=1cm]
            \node[bbox] (N1) {\code{Def}}
              child{node[unknown, xshift=-0.25in] {\code{$\dollar$F}}}
              child{
                node[bbox, fill=bgRed, yshift=-0.1in] {\code{Args}\strut}
                  child{node[unknown2, xshift=-0.4in] {\code{...}}}
                  child{node[bbox, yshift=-0.1in] {\code{DefaultParam}}
                    child{node[unknown] {\code{$\dollar$ARG}}}
                    child{node[hex] {\code{[]}}}
                  }
                  child{node[unknown2, xshift=0.4in] {\code{...}}}
              }
              child{node[unknown2, xshift=0.25in] {\code{...}}}
            ;
          \end{tikzpicture}
       \end{minipage}
       \begin{minipage}{0.52\textwidth}
         \centering
          \begin{tikzpicture}[level distance=1cm]
            \node[bbox] (N1) {\code{Def}}
              child{node[hex, xshift=-0.4in, yshift=0.2in] {\code{append_to}}}
              child{
                node[bbox, fill=bgRed, yshift=-0.1in, xshift=-0.2in] {\code{Args}\strut}
                  child{node[hex, xshift=-0.6in] {\code{element}}}
                  child{node[bbox, yshift=-0.4in] {\code{DefaultParam}}
                    child{node[hex] {\code{to}}}
                    child{node[hex] {\code{[]}}}
                  }
              }
              child{node[bbox, xshift=0.25in, yshift=0.2in] {\code{Seq}\strut}
                child{node[bbox, xshift=-0.1in] {\code{Call}}
                  child{node {\textellipsis}}
                }
                child{node[bbox, xshift=0.1in] {\code{Return}}
                  child{node {\textellipsis}}
                }
              }
            ;
         \end{tikzpicture}
       \end{minipage}
     \end{minipage}
     \begin{minipage}[t][0.3\textheight][t]{\textwidth}
       \begin{minipage}{0.44\textwidth}
         \centering
        \begin{pythoncodeblock}
          def $\dollar$F(..., $\dollar$ARG=[], ...):
            ...
          $\quad$

        \end{pythoncodeblock}
       \end{minipage}
       \begin{minipage}{0.52\textwidth}
        \begin{pythoncodeblock}
          def append_to(element, to=[]):
            to.append(element)
            return to
        \end{pythoncodeblock}
       \end{minipage}
     \end{minipage}
   \end{center}
 \end{frame}

 \begin{frame}[fragile]
   \frametitle{A Default Analysis}

   \begin{center}
    %  \begin{minipage}[t][0.21\textheight][t]{\textwidth}
    %    \raggedright
    %    {\color{blue}\href{https://semgrep.dev/playground/s/Rer0K}{But in this case}}, there would be no match, because the {\color{blue}blue} and
    %    {\color{magenta}magenta} subtrees do not match. In particular, \tikz[baseline] \node[hex, scale=0.7, anchor=base] {\code{f}};
    %    does not match \tikz[baseline] \node[hex, scale=0.7, anchor=base] {\code{g}};.

    %  \end{minipage}
     \begin{minipage}[t][0.65\textheight][t]{\textwidth}
       \begin{minipage}{0.44\textwidth}
         \centering
          \begin{tikzpicture}[level distance=1cm]
            \node[bbox] (N1) {\code{Def}}
              child{node[unknown, xshift=-0.25in, highlight2] {\code{$\dollar$F}}}
              child{
                node[bbox, fill=bgRed, yshift=-0.1in] {\code{Args}\strut}
                  child{node[unknown2, xshift=-0.4in, highlight3] {\code{...}}}
                  child{node[bbox, yshift=-0.1in] {\code{DefaultParam}}
                    child{node[unknown, highlight5] {\code{$\dollar$ARG}}}
                    child{node[hex] {\code{[]}}}
                  }
                  child{node[unknown2, xshift=0.4in] {\code{...}}}
              }
              child{node[unknown2, xshift=0.25in, highlight4] {\code{...}}}
            ;
          \end{tikzpicture}
       \end{minipage}
       \begin{minipage}{0.52\textwidth}
         \centering
          \begin{tikzpicture}[level distance=1cm]
            \node[bbox] (N1) {\code{Def}}
              child{node[hex, xshift=-0.4in, yshift=0.2in, highlight2] {\code{append_to}}}
              child{
                node[bbox, fill=bgRed, yshift=-0.1in, xshift=-0.2in] {\code{Args}\strut}
                  child{node[hex, xshift=-0.6in, highlight3] {\code{element}}}
                  child{node[bbox, yshift=-0.4in] {\code{DefaultParam}}
                    child{node[hex, highlight5] {\code{to}}}
                    child{node[hex] {\code{[]}}}
                  }
              }
              child{node[bbox, xshift=0.25in, yshift=0.2in, highlight4] {\code{Seq}\strut}
                child{node[bbox, xshift=-0.1in, highlight4] {\code{Call}}
                  child{node {\textellipsis}}
                }
                child{node[bbox, xshift=0.1in, highlight4] {\code{Return}}
                  child{node {\textellipsis}}
                }
              }
            ;
         \end{tikzpicture}
       \end{minipage}
     \end{minipage}
     \begin{minipage}[t][0.3\textheight][t]{\textwidth}
       \begin{minipage}{0.44\textwidth}
         \centering
        \begin{pythoncodeblock}
          def !$\dollar$F!(/.../, @$\dollar$ARG@=[], ...):
            ?...?
          $\quad$

        \end{pythoncodeblock}
       \end{minipage}
       \begin{minipage}{0.52\textwidth}
        \begin{pythoncodeblock}
          def !append_to!(/element/, @to@=[]):
            ?to.append(element)?
            ?return to?
        \end{pythoncodeblock}
       \end{minipage}
     \end{minipage}
   \end{center}
 \end{frame}


%  \begin{frame}[fragile]
%    \frametitle{A 150 Analysis}

%    \begin{center}
%      \begin{minipage}[t][0.1\textheight][t]{\textwidth}
%        \raggedright

%        {\color{blue}\href{https://semgrep.dev/playground/s/Abje4}{Let's see it in action:}}
%      \end{minipage}
%      \begin{minipage}[t][0.5\textheight][t]{\textwidth}
%        \begin{minipage}{0.48\textwidth}
%          \centering
%          \begin{tikzpicture}
%            \node[circ] (N3) {\code{@}}
%              child{node[unknown, xshift=-0.35in] (N4) {\code{$\dollar$E1}}}
%              child{node[bbox, xshift=0.35in] (N7) {\code{List}}
%                child{node[unknown] (N8) {\code{$\dollar$E2}}}
%              }
%            ;
%          \end{tikzpicture}
%        \end{minipage}
%        \begin{minipage}{0.48\textwidth}
%          \centering
%          \begin{tikzpicture}[level distance=1cm]
%            \node[] (N1) {\textellipsis}
%              child[missing]
%              child{node[circ, highlight] (N3) {\code{@}}
%                child{node[bbox, xshift=-0.35in, highlight] (N4) {\code{FuncCall}}
%                  child{node[hex, highlight] (N5) {\code{rev}}}
%                  child{node[hex, highlight] (N6) {\code{xs}}}
%                }
%                child{node[bbox, xshift=0.35in, highlight] (N7) {\code{List}}
%                  child{node[hex, highlight] (N8) {\code{x}}}
%                }
%              }
%            ;
%          \end{tikzpicture}
%        \end{minipage}
%      \end{minipage}
%      \begin{minipage}[t][0.3\textheight][t]{\textwidth}
%        \begin{minipage}{0.48\textwidth}
%          \centering
%          \begin{pythoncodeblock}
%            $\dollar$E1 @ [$\dollar$E2]
%          \end{pythoncodeblock}
%        \end{minipage}
%        \begin{minipage}{0.48\textwidth}
%          \begin{pythoncodeblock}
%            fun rev [] = []
%              | rev (x::xs) =
%                  `rev xs @ [x]`
%          \end{pythoncodeblock}
%        \end{minipage}
%      \end{minipage}
%    \end{center}
%  \end{frame}

 \begin{frame}[fragile]
   \frametitle{Matching, Generically}

   Because Semgrep represents all programming languages in the same way,
   matching only needs to be written once.

   \pause
   \vspace{\fill}
   Essentially, we have a function of type:
   \begin{codeblock}
     match_exprs(expr, expr) -> finding list
   \end{codeblock}

   which can do matching for Java, for C, for OCaml, for Python, and many
   other languages besides.
 \end{frame}

 \begin{frame}[fragile]
   \frametitle{Semgrep and Semantics}

   When it comes to program analysis, tree matching isn't necessarily enough.
   Semgrep's main matching capability is at its core more of a Type C analysis,
   but sometimes we want to be able to find more semantic bugs!

   \pause
   \vspace{\fill}

   When it comes to tracking the flow of data, Semgrep uses the theory of
   dataflow analysis as developed before to implement \term{taint tracking},
   which allows you to specify certain sources of bad data, and sinks where that
   data should not flow into, during the program's execution.

   {\color{blue}\href{https://semgrep.dev/playground/s/D1wq}{Here's an example of taint tracking with Semgrep.}}

   \pause
   \vspace{\fill}

   It can also perform constant propagation much in the same way that we did
   earlier:
   {\color{blue}\href{https://semgrep.dev/playground/s/4APL}{Here's an example.}}

   \pause
   \vspace{\fill}

   Semgrep is aware of the semantics of the programming language that it scans,
   meaning that it is not just a simple matching process. It's a
   syntactically-based analysis that is supplemented by deeper semantic analysis.
 \end{frame}

 \sectionSlide{6}{Conclusions}

 \begin{frame}[fragile]
   \frametitle{Program Analysis at Large}

   Program analysis is an extremely important field which aims to both
   enforce that the code we write is safe, as well as supplement the process
   of writing code by providing programmers with the information they need to
   execute smoothly. Even in the face of mathematical impossibility, it
   prevails.

   \pause
   \vspace{\fill}

   Syntax highlighting, IDE warnings, type-checkers, and autoformatters are all
   among the niceties that programmers enjoy, due to advances made in analyzing
   and better making information available about programs.
 \end{frame}

\thankyou

\end{document}
